\documentclass[11pt,a4paper]{article}

\usepackage[utf8]{inputenc}
\usepackage[T1]{fontenc}
\usepackage{amsmath,amssymb,amsthm}
\usepackage{mathtools}
\usepackage{bm}
\usepackage{booktabs}
\usepackage{longtable}
\usepackage{hyperref}
\usepackage[margin=1in]{geometry}
\usepackage{enumitem}
\usepackage{natbib}
\usepackage{xcolor}

\newcommand{\vect}[1]{\bm{#1}}
\newcommand{\mat}[1]{\mathbf{#1}}
\newcommand{\RR}{\mathbb{R}}
\newcommand{\norm}[1]{\lVert #1 \rVert}
\newcommand{\abs}[1]{\lvert #1 \rvert}
\newcommand{\dd}{\mathrm{d}}

\theoremstyle{definition}
\newtheorem{remark}{Remark}

\title{Literature Review and Experimental Synthesis:\\
  Newton--Raphson Minimization on the DFTB0 PES\\
  from Displaced Transition-State Geometries}
\author{M.\ Dincer}
\date{March 2026 --- v12b}

\begin{document}
\maketitle

\tableofcontents
\newpage

%=============================================================================
\section{Introduction}
\label{sec:intro}
%=============================================================================

This document synthesizes insights from two rounds of literature review
(approximately 20 papers) and maps them explicitly to our experimental
findings from versions~1 through~12b of the minimization code.
The goal is to create a permanent, self-contained record of:
\begin{enumerate}[nosep]
  \item what the literature says about our problem class,
  \item what we implemented and why,
  \item what we explicitly chose \emph{not} to implement and why,
  \item what remains open.
\end{enumerate}

\subsection{Problem Statement}

Given a molecule at a geometry displaced by $d \in \{0.5, 1.0, 1.5, 2.0\}$\,\AA\
from a known energy minimum, find the nearest local minimum on the DFTB0
potential energy surface (PES).
\begin{itemize}[nosep]
  \item \textbf{Dataset}: Transition1x~\citep{transition1x} --- 10,073 organic
    reactions, $\leq 7$ heavy atoms, NEB reaction paths computed at
    $\omega$B97x/6-31G(d).
  \item \textbf{Calculator}: DFTB0 via SCINE/Sparrow (analytical Hessian, no
    configuration).
  \item \textbf{Optimizer}: Newton--Raphson in Cartesian coordinates with full
    Hessian at every step, Eckart projection to $3N{-}6$ vibrational modes.
  \item \textbf{Convergence}: $\norm{\vect{f}}_\infty < 10^{-4}$\,eV/\AA\
    AND either (strict) $n_\text{neg} = 0$ or (relaxed) $\lambda_\text{min}
    \geq -\tau_\text{relax}$.
  \item \textbf{Scale}: 287--300 samples per configuration, 10\,000 maximum
    steps.
\end{itemize}

The central experimental finding, established across versions~6--9 (Section~\ref{sec:experiments}), is:
\begin{center}
\fbox{\parbox{0.85\textwidth}{%
  \textbf{Strict convergence} ($n_\text{neg}=0$) with 10k steps drops from
  $\sim$89\% at 0.5\,\AA\ to $\sim$44\% at 2.0\,\AA.
  With 20k steps, convergence reaches \textbf{95\%/85\%/77\%/64\%}
  at 0.5/1.0/1.5/2.0\,\AA\ (v11, RFO+Schlegel+PLS).
  With 50k steps (v12), convergence reaches
  \textbf{95\%/96\%/93\%} at 0.5/1.0/1.5\,\AA\ (RFO+PLS).
  At 2.0\,\AA\ with 20k steps, adaptive probe kicks with late escape
  (v12b) reach \textbf{76\%} (218/287).
  \textbf{Relaxed convergence} ($\lambda_\text{min} \geq -0.01$) is
  $\sim$90\% at \emph{all} noise levels.
  The performance cliff is a \emph{criterion problem}, not an optimizer
  problem.%
}}
\end{center}

Furthermore, the user has demonstrated that 100\% convergence is achievable
via brute-force ``v2-kicking'' (repeatedly perturbing and re-optimizing),
confirming the optimizer can reach valid minima given enough restarts.

%=============================================================================
\section{Theoretical Framework}
\label{sec:theory}
%=============================================================================

\subsection{Newton--Raphson as Spectral Decomposition}
\label{sec:nr-spectral}

Given energy $E(\vect{x})$, gradient $\vect{g} = \nabla E$, and Hessian
$\mat{H} = \nabla^2 E$, the Newton step is
\begin{equation}
  \Delta\vect{x} = -\mat{H}^{-1}\vect{g}.
\end{equation}
After Eckart projection to the vibrational subspace ($3N{-}6$ modes), we
decompose $\mat{H}_\text{vib} = \mat{V}\mat{\Lambda}\mat{V}^\top$ with
eigenvalues $\{\lambda_i\}$ and eigenvectors $\{\vect{v}_i\}$:
\begin{equation}
  \Delta\vect{x} = -\sum_{i=1}^{3N-6} w_i \, (\vect{g}\cdot\vect{v}_i) \, \vect{v}_i,
  \qquad w_i = \frac{1}{\lambda_i}.
  \label{eq:spectral-step}
\end{equation}
When $\mat{H}$ has negative eigenvalues (which occurs $\sim$100\% of the time at our
displacements, with median $n_\text{neg} \approx 12$), the pure Newton step
ascends along negative-eigenvalue directions.  All step builders below modify
the weights $w_i$ to force descent.

\subsection{Our Method = GQT (Goldfeld--Quandt--Trotter)}
\label{sec:gqt}

Our ``shifted Newton'' step (v3 onwards, best performer) sets:
\begin{equation}
  w_i = \frac{1}{\lambda_i + \sigma}, \qquad
  \sigma = \max(0, -\lambda_\text{min}) + \varepsilon,
  \label{eq:shifted-newton}
\end{equation}
with $\varepsilon = 10^{-3}$ (empirically optimal).
This is the spectral form of the \textbf{GQT method}
(Goldfeld--Quandt--Trotter, 1966), which solves the trust-region subproblem
\begin{equation}
  \min_{\vect{s}} \; \vect{g}^\top\vect{s} + \tfrac{1}{2}\vect{s}^\top\mat{H}\vect{s}
  \quad\text{s.t.}\quad \norm{\vect{s}} \leq \Delta.
\end{equation}
The shift $\sigma$ is the Lagrange multiplier.  Our fixed
$\varepsilon$ approximates the optimal $\sigma$ for a particular trust radius.

\textbf{Key property}: Negative modes get \emph{larger} weights than
equal-magnitude positive modes.  Example: $\lambda = -0.01$,
$\sigma = 0.011$ $\Rightarrow$ $w = 1000$; $\lambda = +0.01$,
$\sigma = 0.011$ $\Rightarrow$ $w \approx 48$.  This amplification of
negative modes drives the optimizer toward flipping them---exactly the right
behaviour for escaping saddle regions.

\textbf{Complexity}: GQT has worst-case complexity
$O(\varepsilon^{-5/2})$~\citep{cgt-optimal}, suboptimal compared to ARC's
$O(\varepsilon^{-3/2})$.

\subsection{ARC (Adaptive Regularization with Cubics)}
\label{sec:arc}

ARC~\citep{arc-part1} replaces both the trust region and the shift with a
cubic regularization term:
\begin{equation}
  \min_{\vect{s}} \; m_k(\vect{s}) = E_k + \vect{g}_k^\top\vect{s}
    + \tfrac{1}{2}\vect{s}^\top\mat{H}_k\vect{s}
    + \tfrac{\sigma_k}{3}\norm{\vect{s}}^3.
  \label{eq:arc-model}
\end{equation}
The minimizer satisfies $(\mat{H}_k + \lambda^*\mat{I})\vect{s}^* = -\vect{g}_k$
where $\lambda^* = \sigma_k\norm{\vect{s}^*}$.  In the spectral basis, the
\textbf{secular equation} becomes:
\begin{equation}
  \sigma \norm{\vect{s}(\lambda)} = \lambda,
  \qquad
  \norm{\vect{s}(\lambda)}^2 = \sum_i \frac{c_i^2}{(\lambda_i + \lambda)^2},
  \label{eq:arc-secular}
\end{equation}
where $c_i = \vect{g}\cdot\vect{v}_i$.  We solve~\eqref{eq:arc-secular}
via safeguarded Newton iteration.

\textbf{$\sigma$-adaptation}: After each step, compute the quality ratio
\begin{equation}
  \rho_k = \frac{E(\vect{x}_k) - E(\vect{x}_k + \vect{s}_k)}{m_k(\vect{0}) - m_k(\vect{s}_k)}.
\end{equation}
Update:
\begin{equation}
  \sigma_{k+1} = \begin{cases}
    \sigma_k \cdot \gamma_2 & \text{if } \rho_k > \eta_2 \text{ (very successful)},\\
    \sigma_k & \text{if } \eta_1 \leq \rho_k \leq \eta_2 \text{ (successful)},\\
    \sigma_k \cdot \gamma_1 & \text{if } \rho_k < \eta_1 \text{ (unsuccessful)}.
  \end{cases}
\end{equation}
Default: $\eta_1=0.1$, $\eta_2=0.9$, $\gamma_1=2.0$, $\gamma_2=0.5$.
This is \textbf{self-regulating}: no trust radius, no floor, no collapse
spiral.

\textbf{Complexity}: $O(\varepsilon^{-3/2})$---\textbf{optimal} among
second-order methods~\citep{cgt-optimal}.

\textbf{Implementation}: \texttt{\_nr\_step\_arc()} in
\texttt{minimization.py}, with the secular equation solved iteratively
(lines~700--780).  The ARC step-size control (lines~2350--2450) replaces the
trust region entirely.

\subsection{RFO (Rational Function Optimization)}
\label{sec:rfo}

RFO~\citep{schlegel2011} uses the augmented Hessian eigenvalue problem:
\begin{equation}
  \begin{pmatrix} \mat{H} & \vect{g} \\ \vect{g}^\top & 0 \end{pmatrix}
  \begin{pmatrix} \vect{s} \\ 1 \end{pmatrix}
  = \mu \begin{pmatrix} \vect{s} \\ 1 \end{pmatrix}.
  \label{eq:rfo-augmented}
\end{equation}
For minimization, take the \textbf{lowest eigenvalue} $\mu$.  In the spectral
basis of $\mat{H}$, this yields the secular equation:
\begin{equation}
  \sum_i \frac{c_i^2}{\lambda_i - \mu} + \mu = 0,
  \qquad \mu < \lambda_\text{min}.
  \label{eq:rfo-secular}
\end{equation}
The step is then:
\begin{equation}
  s_i = \frac{-c_i}{\lambda_i - \mu},
  \label{eq:rfo-step}
\end{equation}
which is equivalent to shifted Newton with a \emph{data-dependent} shift
$\sigma = -\mu$.

\textbf{Advantages over fixed-$\varepsilon$ GQT}:
\begin{itemize}[nosep]
  \item Automatic shift determination---no hyperparameter $\varepsilon$.
  \item The shift adapts to the local Hessian spectrum.
  \item Standard in Gaussian, ORCA, and Sella~\citep{sella}.
\end{itemize}

\textbf{Implementation}: \texttt{\_solve\_rfo\_secular()} and
\texttt{\_nr\_step\_rfo()} in \texttt{minimization.py}.  RFO uses the same
trust-region step-size control as the v9 baseline.

\subsection{LM$\leftrightarrow$Cubic Equivalence}
\label{sec:lm-cubic}

Benson and Shanno~\citep{benson-shanno} prove that any
Levenberg--Marquardt step $(\mat{H}+\lambda\mat{I})\vect{s}=-\vect{g}$ is
equivalent to cubic regularization with $M = 2\lambda/\norm{\vect{s}}$.
Conversely, our fixed $\varepsilon = 10^{-3}$ implies variable cubic
strength depending on the step norm:
\begin{equation}
  M_\text{implicit} = \frac{2\varepsilon}{\norm{\vect{s}}}.
\end{equation}
This means our shifted Newton is \emph{already doing implicit cubic
regularization}, but with a strength that varies inversely with step norm
rather than being controlled via $\rho$.

\subsection{The $\varepsilon^{1/2}$ Second-Order Theorem}
\label{sec:eps-half}

Boumal et al.~\citep{boumal-arc-manifolds} prove that when an ARC-type
method drives $\norm{\vect{g}} \leq \varepsilon$, it simultaneously
guarantees:
\begin{equation}
  \lambda_\text{min}(\mat{H}) \geq -\sqrt{\varepsilon}.
  \label{eq:eps-half}
\end{equation}
With our force convergence threshold $\varepsilon = 10^{-4}$\,eV/\AA:
\begin{equation}
  \sqrt{10^{-4}} = 0.01.
\end{equation}
This is \textbf{exactly our empirically optimal relaxed threshold}.
The cascade data (Table~\ref{tab:cascade}) confirms: the convergence rate
plateaus at $\lambda_\text{min} \geq -0.01$, with $\sim$90\% at all noise
levels.  This is not empirical coincidence---it is the theoretical prediction.

\begin{remark}
  This result provides principled justification for relaxed convergence.
  Demanding $n_\text{neg}=0$ is \emph{more strict} than what the optimization
  theory guarantees.  If the optimizer has driven forces below $10^{-4}$\,eV/\AA,
  eigenvalues in $[-0.01, 0)$ are \emph{expected} and acceptable.
\end{remark}

%=============================================================================
\section{Step-Size Control}
\label{sec:step-control}
%=============================================================================

\subsection{Trust Region}
\label{sec:trust-region}

The trust region caps the step to $\norm{\Delta\vect{x}}_\infty \leq \Delta$
(per-atom maximum displacement, using infinity norm~\citep{sella}).  After the
step, compute:
\begin{equation}
  \rho = \frac{E(\vect{x}) - E(\vect{x}+\vect{s})}{m(\vect{0}) - m(\vect{s})}.
\end{equation}

\subsubsection{Original Rules (v1--v9)}
\begin{align}
  \rho > 0.75 &\implies \Delta \gets \min(1.5\Delta, \Delta_\text{max}), \\
  \rho < 0.25 &\implies \Delta \gets \max(0.5\Delta, \Delta_\text{floor}), \\
  \rho < 0    &\implies \Delta \gets \max(0.25\Delta, \Delta_\text{floor}).
\end{align}
With $\Delta_\text{max} = 1.3$\,\AA\ and $\Delta_\text{floor} = 0.01$\,\AA.

\subsubsection{Schlegel (2011) Corrections}
\label{sec:schlegel-tr}

Schlegel's review~\citep{schlegel2011} identifies two problems with standard
rules:

\paragraph{Growth condition.}
Growing the trust radius when $\rho > 0.75$ but the step is far from the
boundary ($\norm{\vect{s}} \ll \Delta$) is pointless and can cause
instability.  The corrected rule:
\begin{equation}
  \text{Grow only if } \rho > 0.75 \text{ AND }
  \norm{\vect{s}}_\infty \geq 0.8\,\Delta.
  \label{eq:schlegel-grow}
\end{equation}

\paragraph{Shrink formula.}
Shrinking to a fraction of $\Delta$ (the radius) ignores the actual step
taken.  If the step was much smaller than $\Delta$ and still bad, shrinking
$\Delta$ doesn't help.  The corrected rule:
\begin{equation}
  \Delta \gets \max(0.25\,\norm{\vect{s}}_\infty, \, \Delta_\text{floor}).
  \label{eq:schlegel-shrink}
\end{equation}

\textbf{Implementation}: Gated by \texttt{schlegel\_trust\_update=True}
parameter (v10b).  The flag ensures backward compatibility.

\subsection{ARC $\sigma$-Adaptation}
\label{sec:arc-sigma}

ARC replaces the trust region entirely (Section~\ref{sec:arc}).  The
regularization parameter $\sigma$ controls effective step size: large $\sigma$
$\Rightarrow$ short, cautious steps; small $\sigma$ $\Rightarrow$ long,
aggressive steps.  There is no radius floor, so the collapse spiral
(Section~\ref{sec:collapse}) cannot occur.

\subsection{Line Search}
\label{sec:line-search}

Armijo backtracking: $E(\vect{x}+\alpha\vect{d}) \leq E(\vect{x}) + c_1
\alpha\, \vect{g}^\top\vect{d}$ with $c_1=10^{-4}$.

\textbf{v7 result}: Line search is \emph{uniformly worse} than trust region
at all noise levels (Table~\ref{tab:v7}).  Reason: line search commits to the
Newton direction $\vect{d}$, but at large displacement, this direction has
cosine similarity $\sim$0.03--0.11 with the true direction
(Table~\ref{tab:direction-quality}).  Trust region limits damage to $\Delta$;
line search amplifies the wrong direction.

\subsection{Trust Radius Collapse Mechanism}
\label{sec:collapse}

The dominant failure mode (identified in v6 analysis):
\begin{enumerate}[nosep]
  \item Shifted Newton weight $w_i = 1/(\lambda_i+\sigma)$, max = $1/\varepsilon = 1000\times$.
  \item Near-zero eigenvalue modes get amplified $1000\times$, producing oversized steps.
  \item Oversized step violates quadratic model: $\rho < 0.25$ $\Rightarrow$ $\Delta \gets 0.25\Delta$.
  \item After $\sim$4 bad steps: $1.3 \to 0.33 \to 0.08 \to 0.02 \to 0.01$\,\AA\ (floor).
  \item At floor, all modes scaled uniformly. Neg-mode displacement $\approx 5\times10^{-4}$\,\AA/step.
  \item Recovery needs $\sim$13 consecutive good steps. But eigenvalues oscillate $\Rightarrow$ never happens.
\end{enumerate}

This is why ARC and RFO are promising: they control regularization
\emph{continuously} without a hard floor.

%=============================================================================
\section{GDIIS Theory and Implementation}
\label{sec:gdiis}
%=============================================================================

\subsection{Pulay's DIIS for SCF vs.\ Cs\'asz\'ar--Pulay's GDIIS}

\textbf{DIIS} (Direct Inversion in the Iterative Subspace)~\citep{pulay-diis}
was developed for SCF convergence: find coefficients $\{c_i\}$ with
$\sum c_i = 1$ minimizing $\norm{\sum c_i \vect{e}_i}^2$ where $\vect{e}_i$
are error vectors.  The interpolated solution is $\vect{x}_\text{DIIS} =
\sum c_i \vect{x}_i$.

\textbf{GDIIS}~\citep{gdiis-csaszar-pulay} extends this to geometry
optimization.  The key difference is the error vectors.

\subsection{Schlegel's Correct Formulation}

Schlegel~\citep{schlegel2011} specifies (Eq.~16) that the error vectors
for GDIIS should be the \emph{Newton step residuals} $\mat{H}^{-1}\vect{g}_i$
(or equivalently, the Newton step vectors), not the raw gradients:
\begin{equation}
  B_{ij} = (\mat{H}^{-1}\vect{g}_i)^\top (\mat{H}^{-1}\vect{g}_j).
  \label{eq:gdiis-B}
\end{equation}
Solve $\mat{B}\vect{c} = \vect{0}$ subject to $\sum c_i = 1$ via the
constrained system:
\begin{equation}
  \begin{pmatrix} \mat{B} & \vect{1} \\ \vect{1}^\top & 0 \end{pmatrix}
  \begin{pmatrix} \vect{c} \\ \lambda \end{pmatrix}
  = \begin{pmatrix} \vect{0} \\ 1 \end{pmatrix}.
  \label{eq:gdiis-system}
\end{equation}
The interpolated geometry is $\vect{x}_\text{GDIIS} = \sum c_i \vect{x}_i$.

\subsection{Our Bug: Gradient-Based Error Vectors}

Our original \texttt{\_GDIISExtrapolator} (v5/v10 initial implementation)
used \emph{raw gradients} $\vect{g}_i$ as error vectors:
$B_{ij} = \vect{g}_i^\top\vect{g}_j$.  This is Pulay's \emph{SCF-DIIS}
formulation, not geometry-optimization GDIIS.  The difference matters because:
\begin{itemize}[nosep]
  \item Raw gradient dot products ignore the Hessian spectrum: two points
    with parallel gradients but different curvatures get the same B-matrix
    entry.
  \item Newton step dot products incorporate curvature: high-curvature modes
    are de-emphasized, low-curvature modes amplified.
\end{itemize}

\textbf{Fix (v10b)}: Modified \texttt{store()} to accept the Newton step
vector; B-matrix now uses step dot products per~\eqref{eq:gdiis-B}.
Backward-compatible: if step not provided, falls back to gradient.

\subsection{GEDIIS (Li--Frisch, 2006)}
\label{sec:gediis}

GEDIIS~\citep{li-frisch-gediis} is an energy-based DIIS variant that
interpolates geometries to minimize \emph{energy} rather than error norm.
Hybrid GEDIIS/GDIIS switches based on convergence stage.
We did not implement GEDIIS but note it as a possible future extension.

\subsection{Late-Stage Acceleration}
\label{sec:gdiis-late}

Schlegel~\citep{schlegel2011} describes GDIIS as routinely accelerating
\emph{final convergence}, not just breaking oscillations.
Our v10 implementation only triggered GDIIS on oscillation detection.

\textbf{v10b addition}: New parameter \texttt{gdiis\_late\_force\_threshold}
(default 0).  When \texttt{force\_norm < threshold}, GDIIS is attempted every
\texttt{gdiis\_every} steps regardless of oscillation.  This implements
Schlegel's recommendation.

%=============================================================================
\section{Ghost Modes and Convergence Criteria}
\label{sec:ghost-modes}
%=============================================================================

\subsection{Eckart Projection}

Eckart projection removes 6 translation/rotation modes from the $3N$-dimensional
Hessian, yielding $3N{-}6$ vibrational modes.  The projection is
\emph{exact only at a stationary point}.  Away from equilibrium,
rotation--vibration coupling leaks into the vibrational subspace,
creating small spurious eigenvalues.

Implementation: \texttt{get\_vib\_evals\_evecs()} from
\texttt{multi\_mode\_eckartmw.py}, with optional \texttt{purify\_hessian}
symmetrization (found unhelpful in practice).

\subsection{ORCA Acoustic Sum Rule}

The ORCA manual~\citep{orca-manual} explicitly mentions an ``acoustic sum rule''
toggle to reduce noise-induced imaginary frequencies---the same phenomenon.
This confirms ghost modes are a known numerical artifact, not a deficiency of
our code.

\subsection{Morse--Bott Theory}
\label{sec:morse-bott}

Luo et al.~\citep{convergence-degenerate} analyze convergence near degenerate
critical points (eigenvalue zero) using Morse--Bott theory.  Key result:
\begin{quote}
  Near-zero eigenvalues on critical manifolds are \emph{neutral} for
  convergence.  The gradient component in zero-eigenvalue directions is
  $O(r^2)$ where $r$ is distance to the manifold.
\end{quote}
This directly explains why ghost eigenvalues in $[-10^{-4}, 0)$ do not
prevent force convergence.  The optimizer converges in the non-degenerate
(vibrational) subspace; the degenerate (ghost) subspace produces tiny,
ignorable gradient contributions.

\subsection{The $n_\text{neg}=0$ Criterion is Inappropriate}
\label{sec:nneg-inappropriate}

Three lines of evidence converge:
\begin{enumerate}[nosep]
  \item \textbf{Theoretical}: The $\varepsilon^{1/2}$ theorem
    (Section~\ref{sec:eps-half}) predicts eigenvalues in $[-0.01, 0)$ are
    acceptable when $\norm{\vect{g}} \leq 10^{-4}$.
  \item \textbf{Numerical}: Ghost modes from Eckart projection artifacts
    populate the $[-10^{-4}, 0)$ band.
  \item \textbf{Cross-surface}: At DFT-labeled minima, 90\% have $n_\text{neg}>0$
    under DFTB0, with \emph{strongly negative} eigenvalues ($<-0.01$).  These
    are genuine PES disagreements between theory levels, not optimizer failures.
\end{enumerate}

The cascade evaluation (Table~\ref{tab:cascade}) quantifies the effect:
the gap between strict ($n_\text{neg}=0$) and relaxed ($\lambda_\text{min}
\geq -0.01$) grows from 18 percentage points at 0.5\,\AA\ to 73 points at
2.0\,\AA.

\subsection{Saddle Point Convergence ($n_\text{neg}=1$)}

For mode-following, GAD (Gentlest Ascent Dynamics), or saddle-point search
methods, the convergence criterion would be $n_\text{neg}=1$ (first-order
saddle point).  Our codebase is currently minimization-only, but the
classification infrastructure (\texttt{cascade\_at\_convergence}) already
supports arbitrary eigenvalue counting.

%=============================================================================
\section{Literature Confirmations of Experimental Findings}
\label{sec:confirmations}
%=============================================================================

\begin{longtable}{p{0.35\textwidth}p{0.6\textwidth}}
\toprule
\textbf{Our Result} & \textbf{Literature Explanation} \\
\midrule
\endhead

iHiSD crossover hurts (v8: pure SN 43.3\% $\to$ cm=2.0 only 15.9\%)
& Su et al.~\citep{ihisd}: iHiSD assumes gradient alignment tendency near
  saddle points.  Our data shows cosine similarity 0.03--0.11 at
  0.5--2.0\,\AA\ displacement (Table~\ref{tab:direction-quality}).
  We are outside the theory's applicability regime. \\

First-order methods fail (PIC-ARC: 0\%)
& Packwood~\citep{packwood2018}: even approximate curvature dramatically
  outperforms gradients alone in nonconvex landscapes.
  Cartis et al.~\citep{cgt-complexity}: Newton can be as slow as steepest
  descent in worst case, but SD can never be faster than regularized Newton. \\

SPDN weight cap catastrophic (v5: 0\%)
& Capping weights at $1/\tau_\text{hard} = 100$ discards curvature
  information.  Packwood shows the exact Hessian can be \emph{worse} than a
  ``cleaned-up'' version---but cleaning up means enforcing positive
  definiteness, not capping.  Different operation entirely. \\

LM damping inferior (v2)
& LM shifts \emph{all} eigenvalues uniformly by $\mu$:
  $w_i = |\lambda_i|/(\lambda_i^2+\mu^2)$.  Shifted Newton targets
  only the problematic modes via $\sigma = \max(0,-\lambda_\text{min})+\varepsilon$.
  LM over-regularizes positive eigenvalues. \\

Line search $\prec$ trust region (v7)
& Standard theory: line search requires a descent direction.  Our Newton
  direction is poor at large displacement.  TR limits damage; LS commits to
  the direction~\citep{schlegel2011,cgt-optimal}. \\

\bottomrule
\end{longtable}

%=============================================================================
\section{Dataset and Calculator Considerations}
\label{sec:dataset}
%=============================================================================

\subsection{Transition1x Dataset}
\label{sec:transition1x}

Transition1x~\citep{transition1x} provides 10,073 organic reactions:
\begin{itemize}[nosep]
  \item NEB paths at $\omega$B97X-D3/def2-TZVP (some at $\omega$B97x/6-31G(d)).
  \item Endpoints labeled as ``reactant'' and ``product'' minima.
  \item Transition states verified by single imaginary frequency.
  \item Maximum 7 heavy atoms, 23 total atoms.
\end{itemize}

\textbf{Critical insight}: The dataset performs \emph{no frequency checks on
minima}---only force convergence to 0.01\,eV/\AA.  This means:
\begin{enumerate}[nosep]
  \item Some ``minima'' may be saddle points at the original DFT level.
  \item The 0.01\,eV/\AA\ force threshold is $100\times$ our convergence
    criterion.
  \item These endpoints are then used as references for evaluating our
    DFTB0 optimizer.
\end{enumerate}
Approximately 16\% of NEB reactions failed and were discarded from the
dataset.

\subsection{Level-of-Theory Mismatch}

Three levels of theory are involved:
\begin{center}
  $\omega$B97X-D3/def2-TZVP \quad$\xrightarrow{\text{NEB}}$\quad
  $\omega$B97x/6-31G(d) \quad$\xrightarrow{\text{our optimizer}}$\quad
  DFTB0
\end{center}
Each transition changes the PES.  A minimum at one level may not be a minimum
at another.  Our forensic analysis (Section~\ref{sec:dftb0-forensics})
quantifies this.

\subsection{DFTB0 Surface Forensics}
\label{sec:dftb0-forensics}

At DFT-labeled reactant minima, under DFTB0:
\begin{center}
\begin{tabular}{ll}
\toprule
Metric & Value \\
\midrule
Fraction with $n_\text{neg}=0$ & \textbf{10\%} \\
Median $n_\text{neg}$ & 2 \\
Median force at ``minimum'' & 1.39\,eV/\AA \\
Negative eigenvalue band $\lambda < -0.1$ & 0.90 modes/molecule \\
Negative eigenvalue band $-0.1 \leq \lambda < -0.01$ & 0.63 modes/molecule \\
Ghost band $\lambda \in [-10^{-3}, 0)$ & 0.08 modes/molecule \\
Condition number $\kappa$ & $\sim 10^{3.2}$ \\
\bottomrule
\end{tabular}
\end{center}

The strongly negative eigenvalues ($<-0.01$) confirm genuine PES disagreements
between DFT and DFTB0---these are \emph{not} numerical ghosts.  The ghost band
is nearly empty at actual minima but populated at displaced geometries.

\subsection{SCINE/Sparrow Calculator}

SCINE/Sparrow provides:
\begin{itemize}[nosep]
  \item Zero-configuration DFTB0 (no parameters to set).
  \item Analytical Hessian (not finite difference).
  \item Hessian noise floor: estimated $\sim 8\times10^{-3}$ from the
    \texttt{nr\_threshold} parameter.
\end{itemize}

We did \emph{not} explore SCINE SCF convergence settings, which could
potentially reduce the Hessian noise floor.  This is noted as future work.

%=============================================================================
\section{Experimental Results Summary}
\label{sec:experiments}
%=============================================================================

\subsection{Version History}

\begin{longtable}{clp{5cm}p{3cm}l}
\toprule
\textbf{Version} & \textbf{Key Idea} & \textbf{Result} & \textbf{Verdict} \\
\midrule
\endhead

v1 & Hard-filter NR + TR & Baseline & --- \\
v2 & LM damping + cascade & LM inferior to SN & Diagnostics kept \\
v3 & \textbf{Shifted Newton} + escape & 80\% at 1.0\,\AA\ (15 samples) & \textbf{Best step} \\
v4 & 5 targeted failure fixes & All hurt (53$\to$10\%) & All reverted \\
v5 & SPDN (partition + GDIIS + LS) & 0\% convergence & Catastrophic \\
v6 & Back to basics: SN + TR & 76--79\% at 0.5\,\AA\ (300 samples) & Confirmed \\
v7 & Line search vs.\ trust region & LS worse at all noise & TR wins \\
v8 & iHiSD crossover (GD$\to$NR) & Crossover hurts & Pure SN best \\
PIC-ARC & First-order + cubic reg & 0\% convergence & First-order dead \\
v9 & Relaxed convergence & 80/60/46/29\% strict & Criterion problem \\
v10 & ARC + GDIIS & ARC catastrophic (23/2/0/0\%) & \textbf{ARC dead} \\
v10.2 & +RFO, +Schlegel TR, +GDIIS-late & RFO 67/48/32/12\%; GDIIS zero effect & GDIIS dead \\
v10c & +Polynomial line search & \textbf{89/78/63/42\%} (RFO+Schlegel+PLS) & \textbf{PLS breakthrough} \\
v11 & +20k steps, trust floor, GQT+Schl+PLS & \textbf{95/85/77/64\%} (RFO+Schl+PLS, 20k) & \textbf{20k breakthrough} \\
v12 & +50k steps, osc/blind kicks & \textbf{95/96/93\%} (RFO+PLS, 50k); kicks $+$5.9\,pts at 2.0\,\AA\ only & \textbf{50k breakthrough} \\
v12b & +adaptive kick scaling, probe kicks, late escape & \textbf{76.0\%} at 2.0\,\AA\ (20k); probe is the key ($+$9\,pts) & \textbf{Probe breakthrough} \\
\bottomrule
\end{longtable}

\subsection{Best Convergence Rates (v11)}
\label{sec:best-rates}

\begin{table}[h]
\centering
\begin{tabular}{cccccc}
\toprule
Noise & v9 GQT & v10c Best (10k) & v11 Best (10k) & v11 Best (20k) & $\Delta$ (v9$\to$v11) \\
\midrule
0.5\,\AA & 80\% & 89\% & 89\% & \textbf{95\%} & +15\,pts \\
1.0\,\AA & 60\% & 78\% & 80\% & \textbf{85\%} & +25\,pts \\
1.5\,\AA & 46\% & 63\% & 61\% & \textbf{77\%} & +31\,pts \\
2.0\,\AA & 29\% & 45\% & 44\% & \textbf{64\%} & +35\,pts \\
\bottomrule
\end{tabular}
\caption{Best convergence rates across all versions (287 samples, strict $n_\text{neg}=0$).
  v10c Best (10k) = RFO+Schlegel+PLS at 0.5--1.5\,\AA, GQT+PLS at 2.0\,\AA.
  v11 Best (10k) = RFO+PLS (no Schlegel) at 1.0/2.0\,\AA, RFO+Schlegel+PLS at 0.5/1.5\,\AA.
  v11 Best (20k) = RFO+Schlegel+PLS at all noise levels.
  Note: v10c and v11 use different random seeds; 10k rates agree within $\pm$3.5\,pts.}
\label{tab:best-rates}
\end{table}

\subsection{v10/v10c Comprehensive Results}
\label{sec:v10-results}

Four experimental runs were conducted (287 samples each, 10k steps, DFTB0):

\paragraph{v9 (re-run, 16 configs).}
Baseline shifted Newton with $\varepsilon=10^{-3}$, trust region,
relaxed eigenvalue thresholds varied.
Convergence: 80/60/46/29\% at 0.5/1.0/1.5/2.0\,\AA.
\textbf{All relaxed thresholds} ($-10^{-3}$, $-5\times10^{-3}$, $-10^{-2}$)
\textbf{give identical counts}---the relaxed criterion does not separate
any additional population.

\paragraph{v10 (28 configs).}
First ARC test with $\sigma_0 \in \{0.1, 1.0, 10.0\}$,
$\gamma_1 \in \{2.0, 3.0\}$, GDIIS buffer $\in \{8, 12\}$.
GQT baseline: 82/55/47/37\%.
\textbf{ARC catastrophic}: best 23/2/0/0\% at $n=0.5$.
Failure autopsy: \textbf{80.3\% of all trajectories classified as
\texttt{energy\_plateau}}---$\sigma$ over-regularizes, steps become
too small, energy stalls.

\paragraph{v10.2 (48 configs).}
Added RFO, Schlegel trust radius fixes, GDIIS late-stage trigger.
GQT: 75/53/44/27\%.
RFO: 67/48/32/12\%.
Schlegel TR: mixed ($\pm$2--5\,pts depending on noise level).
\textbf{GDIIS zero effect}: RFO+Schlegel+GDIIS = RFO+Schlegel exactly
(same convergence count, same oscillating failure count).

\paragraph{v10c (56 configs).}
Added polynomial line search (PLS) to GQT and RFO configs.
\textbf{RFO+Schlegel+PLS: 89/78/63/42\%} (new champion at 0.5--1.5\,\AA).
GQT+PLS: 83/67/56/45\% (wins at 2.0\,\AA).

\begin{table}[h]
\centering
\caption{Full convergence comparison across v10c configurations (287 samples).}
\label{tab:v10c-full}
\begin{tabular}{lcccc}
\toprule
Configuration & 0.5\,\AA & 1.0\,\AA & 1.5\,\AA & 2.0\,\AA \\
\midrule
GQT (v9 baseline) & 76\% & 55\% & 49\% & 27\% \\
GQT + Schlegel TR & 78\% & 57\% & 47\% & 29\% \\
GQT + PLS & \textbf{83\%} & \textbf{67\%} & \textbf{56\%} & \textbf{45\%} \\
GQT + Schlegel TR + PLS & 83\% & 68\% & 54\% & 42\% \\ % (v11)
\midrule
RFO & 65\% & 46\% & 32\% & 13\% \\
RFO + Schlegel TR & 67\% & 44\% & 31\% & 14\% \\
RFO + Schlegel TR + PLS & \textbf{89\%} & \textbf{78\%} & \textbf{63\%} & \textbf{42\%} \\
RFO + Schlegel TR + GDIIS & 67\% & 44\% & 31\% & 14\% \\
\midrule
ARC ($\sigma_0=1$, $\gamma_1=2$) & 18\% & 2\% & 0\% & 0\% \\
ARC ($\sigma_0=1$, $\gamma_1=3$) & 20\% & 2\% & 0\% & 0\% \\
\bottomrule
\end{tabular}
\end{table}

\subsubsection{Factor Analysis}

\paragraph{Polynomial line search: the breakthrough.}
PLS is the single most impactful intervention across all versions:
\begin{itemize}[nosep]
  \item GQT $\to$ GQT+PLS: +7/+12/+7/+18 points
  \item RFO+Schlegel $\to$ RFO+Schlegel+PLS: +22/+34/+32/+28 points
  \item PLS halves oscillating failures (autopsy): e.g., at $n=0.5$,
    GQT oscillating 32$\to$12; RFO oscillating 61$\to$18.
  \item Mechanism: cubic interpolation using energy+gradient at current
    and previous geometry; zero extra evaluations; stabilizes
    eigenvalue-crossing regime.
\end{itemize}

\paragraph{ARC: catastrophic failure.}
Across all $\sigma_0$, $\gamma_1$ combinations:
\begin{itemize}[nosep]
  \item Energy plateau is the dominant failure mode (80\% of all ARC
    trajectories in v10).
  \item $\sigma$ grows without bound, steps shrink to near-zero,
    energy stalls.
  \item No $\sigma_0$/$\gamma_1$ combination recovers useful performance.
  \item GDIIS does not rescue ARC (same counts with/without).
\end{itemize}
Root cause: ARC's $\sigma$-adaptation formula was designed for smooth
optimization problems.  On the DFTB0 PES with Hessian noise, even good
steps can have poor $\rho$ values, causing $\sigma$ to ratchet upward.

\paragraph{GDIIS: zero effect.}
Across all configurations (buffer 8, 12; oscillation-triggered;
late-stage threshold 0.01):
\begin{itemize}[nosep]
  \item Convergence counts identical with/without GDIIS.
  \item Failure mode counts identical with/without GDIIS.
  \item Neither oscillation detection nor late-stage trigger produces
    any benefit.
\end{itemize}
Hypothesis: GDIIS extrapolation requires the error vectors to lie in a
low-dimensional subspace.  In our high-dimensional problem ($\sim$60
DOF), the oscillation involves too many coupled modes for 8--12
history vectors to capture.

\paragraph{Schlegel trust radius: mixed.}
Boundary-check growth + step-anchored shrink produces $\pm$2--5\,pts
depending on noise level.  Not a consistent improvement on its own,
but combines well with PLS (RFO+Schlegel+PLS is champion).

\paragraph{RFO vs.\ GQT: crossover at high noise.}
RFO+PLS wins at 0.5--1.5\,\AA\ but GQT+PLS wins at 2.0\,\AA\ (45\%
vs.\ 42\%).  At large displacements, RFO's augmented Hessian shift may
be too aggressive, while GQT's fixed $\varepsilon$ provides a more
conservative regularization floor.

\subsubsection{Failure Autopsy Summary (v10c)}

\begin{table}[h]
\centering
\caption{Failure classification across v10c (16,023 trajectories).}
\label{tab:autopsy-v10c}
\begin{tabular}{lrr}
\toprule
Classification & Count & Fraction \\
\midrule
converged & 4,446 & 27.7\% \\
energy\_plateau (ARC) & 7,595 & 47.4\% \\
oscillating & 2,710 & 16.9\% \\
almost\_converged & 1,000 & 6.2\% \\
ghost\_modes & 241 & 1.5\% \\
slow\_convergence & 26 & 0.2\% \\
drifting & 5 & 0.0\% \\
\bottomrule
\end{tabular}
\end{table}

Excluding ARC configs, the remaining failure landscape is:
oscillating ($\sim$50\% of non-ARC failures),
almost\_converged ($\sim$30\%),
ghost\_modes ($\sim$10\%).
PLS halves the oscillating population; the almost\_converged and
ghost\_mode populations are unaffected by any v10 intervention.

\subsection{Transition1x Frequency Verification (v10c)}
\label{sec:transition1x-verification}

DFTB0 frequency analysis on all 287 DFT-labeled endpoints:
\begin{itemize}[nosep]
  \item Reactants: 29/287 (10.1\%) are genuine minima ($n_\text{neg}=0$).
  \item Products: 30/287 (10.5\%) are genuine minima ($n_\text{neg}=0$).
  \item Correlation with optimizer failure: Pearson $r=0.22$ (weak).
\end{itemize}
The weak correlation confirms that optimizer failure is not primarily
driven by the starting-point PES character.

\subsection{Direction Quality}
\label{sec:direction-quality}

\begin{table}[h]
\centering
\begin{tabular}{ccccc}
\toprule
Noise & $\cos(\vect{g}_\text{SD}, \vect{d}_\text{true})$ & $\cos(\vect{d}_\text{NR}, \vect{d}_\text{true})$ & SD better \% & Starting $n_\text{neg}$ \\
\midrule
0.5\,\AA & 0.029 & 0.055 & 46\% & 12.1 \\
1.0\,\AA & 0.038 & 0.025 & 54\% & 12.4 \\
1.5\,\AA & 0.093 & 0.043 & 62\% & 12.1 \\
2.0\,\AA & 0.110 & 0.027 & \textbf{70\%} & 12.0 \\
\bottomrule
\end{tabular}
\caption{Direction quality at noisy starting points.  Both directions are
  near-orthogonal to truth.  SD ``improves'' with noise but is still too poor
  for useful steps.}
\label{tab:direction-quality}
\end{table}

\subsection{Condition Number Scaling}

\begin{table}[h]
\centering
\begin{tabular}{cccc}
\toprule
Noise & Median $\kappa$ & $\log_{10}(\kappa)$ & Interpretation \\
\midrule
0.5\,\AA & $1.14\times10^4$ & 4.1 & Moderate \\
1.0\,\AA & $2.79\times10^4$ & 4.4 & Degraded \\
1.5\,\AA & $2.00\times10^5$ & 5.3 & Severe \\
2.0\,\AA & $2.30\times10^6$ & 6.4 & Catastrophic \\
True min & $1.55\times10^3$ & 3.2 & Well-conditioned \\
\bottomrule
\end{tabular}
\caption{Hessian condition number vs.\ displacement.  The sharpest convergence
  cliff (59\%$\to$41\%) coincides with $\kappa$ jumping
  $10^{4.4}\to10^{5.3}$.}
\label{tab:condition}
\end{table}

\subsection{Failure Populations}

Three distinct failure populations identified from trajectory analysis:

\paragraph{Population A: ``Blind Modes'' (49.6\% of pre-v10c failures).}
$\lambda_\text{min} \in [-10^{-3}, -2\times10^{-4}]$, forces 0.03--0.09\,eV/\AA,
gradient overlap with negative modes $<0.1$.  Trust radius at floor.
The optimizer ``sees'' the negative eigenvalue but the gradient has no
component in that direction---it cannot flip the mode.
\emph{v10c update}: Corresponds to the \texttt{almost\_converged} and
\texttt{ghost\_modes} populations in the autopsy.
\textbf{Unaffected by PLS, GDIIS, or any v10 intervention.}

\paragraph{Population B: ``Oscillation--Collapse'' (44.2\% of pre-v10c failures).}
$\lambda_\text{min} \in [-10^{-2}, -10^{-3}]$, forces 0.1--0.37\,eV/\AA.
Period-$\sim$8 oscillation cycle detected by FFT.
\emph{v10c update}: \textbf{PLS halves this population} (e.g., GQT at
$n=0.5$: 32$\to$12 oscillating; RFO at $n=0.5$: 61$\to$18).
GDIIS had zero effect despite being designed for exactly this case.

\paragraph{Population C: ``Stagnation Zombies'' (subset).}
Stagnation $>50\%$ of trajectory.  Example: sample~098 frozen 99.6\% of steps.

\paragraph{Population D: ``Energy Plateau'' (ARC-specific).}
$\sigma$ over-regularization causes near-zero steps; energy stalls.
80\% of all ARC trajectories.  This population does not exist for
GQT or RFO configurations.

%=============================================================================
\subsection{v11 Grid Results (mar8 Run)}
\label{sec:v11}
%=============================================================================

\subsubsection{Experimental Design}

The v11 grid tested 8 algorithms $\times$ 4 noise levels = 32 configurations,
each with 287 samples, 10k or 20k maximum steps.  \textbf{All v11 configs use
PLS.}  The grid addresses four questions:
\begin{enumerate}[nosep]
  \item Does GQT+Schlegel+PLS (never tested) outperform GQT+PLS or RFO+Schlegel+PLS?
  \item Does doubling the step budget (20k) rescue almost\_converged or oscillating failures?
  \item Does lowering the trust radius floor from 0.01\,\AA\ to 0.005\,\AA\ help?
  \item Is Schlegel TR helpful when PLS is already active?
\end{enumerate}
Additionally, RFO+PLS (no Schlegel) was tested for the first time.

\subsubsection{Reproducibility Assessment}
\label{sec:reproducibility}

The v11 run used a different random seed than the v10c run.  Two PLS
configs (RFO+Schlegel+PLS and GQT+PLS) appear in both runs, enabling
cross-seed comparison:

\begin{table}[h]
\centering
\caption{Cross-seed reproducibility: v10c (SLURM 1104386) vs.\ v11 (SLURM 1113757)
  for identical configurations.  $N=287$ samples each.}
\label{tab:reproducibility}
\begin{tabular}{lccccc}
\toprule
Config & Noise & v10c & v11 & $\Delta$ (pts) \\
\midrule
RFO+Schl+PLS & 0.5\,\AA & 254/287 (88.5\%) & 256/287 (89.2\%) & $+$0.7 \\
              & 1.0\,\AA & 223/287 (77.7\%) & 225/287 (78.4\%) & $+$0.7 \\
              & 1.5\,\AA & 182/287 (63.4\%) & 175/287 (61.0\%) & $-$2.4 \\
              & 2.0\,\AA & 120/287 (41.8\%) & 124/287 (43.2\%) & $+$1.4 \\
\midrule
GQT+PLS      & 0.5\,\AA & 239/287 (83.3\%) & 242/287 (84.3\%) & $+$1.0 \\
              & 1.0\,\AA & 191/287 (66.6\%) & 191/287 (66.6\%) & $\phantom{+}$0.0 \\
              & 1.5\,\AA & 161/287 (56.1\%) & 156/287 (54.4\%) & $-$1.7 \\
              & 2.0\,\AA & 128/287 (44.6\%) & 118/287 (41.1\%) & $-$3.5 \\
\bottomrule
\end{tabular}
\end{table}

Mean absolute difference: 1.4\,pts.  Maximum: 3.5\,pts at $n=2.0$ for GQT+PLS.
This establishes the \textbf{seed-dependent variation at $\pm$1--4 pts},
growing with noise level.  Rankings are preserved across seeds at 0.5--1.5\,\AA.
At 2.0\,\AA, the GQT~vs.~RFO ordering can flip between seeds (the two methods
are statistically indistinguishable at this noise level with $N=287$).

The failure \emph{distributions} also reproduce across seeds.  For the two
matching configs (RFO+Schlegel+PLS, GQT+PLS), the per-noise-level autopsy:

\begin{table}[h]
\centering
\caption{Cross-seed failure autopsy: v10c vs.\ v11 for matching configs.
  ``Osc'' = oscillating, ``AC'' = almost\_converged, ``GM'' = ghost\_modes.
  Total $N$ differs slightly due to SCINE errors (v10c: 285--287; v11: 285--287).}
\label{tab:autopsy-cross-seed}
\begin{tabular}{llrrrrrrrr}
\toprule
Config & Noise & \multicolumn{2}{c}{Conv} & \multicolumn{2}{c}{Osc} & \multicolumn{2}{c}{AC} & \multicolumn{2}{c}{GM} \\
       &       & v10c & v11 & v10c & v11 & v10c & v11 & v10c & v11 \\
\midrule
RFO+Schl+PLS & 0.5 & 254 & 256 & 18 & 18 & 10 & 9 & 3 & 3 \\
              & 1.0 & 223 & 225 & 40 & 35 & 22 & 21 & 1 & 5 \\
              & 1.5 & 182 & 175 & 76 & 76 & 19 & 23 & 5 & 9 \\
              & 2.0 & 120 & 124 & 121 & 117 & 21 & 14 & 18 & 17 \\
\midrule
GQT+PLS      & 0.5 & 239 & 242 & 12 & 11 & 25 & 25 & 3 & 7 \\
              & 1.0 & 191 & 191 & 32 & 43 & 52 & 46 & 10 & 7 \\
              & 1.5 & 161 & 156 & 76 & 75 & 41 & 45 & 8 & 10 \\
              & 2.0 & 128 & 118 & 98 & 116 & 37 & 26 & 23 & 24 \\
\bottomrule
\end{tabular}
\end{table}

The failure mode counts are remarkably stable across seeds.  The oscillating
count for RFO+Schl+PLS at $n=0.5$ is identical (18 vs.\ 18); at $n=1.5$
it is also identical (76 vs.\ 76).  The largest discrepancy is GQT+PLS
oscillating at $n=2.0$ (98 vs.\ 116), which tracks the convergence
difference ($-$10 samples).  This confirms that the failure classification
is \textbf{reproducible across random seeds}, not an artifact of a
particular displacement realization.

\subsubsection{v11 Convergence Table}
\label{sec:v11-table}

\begin{table}[h]
\centering
\caption{v11 convergence results (287 samples, strict $n_\text{neg}=0$, all
  configs use PLS).  ``trf'' = trust radius floor 0.005\,\AA\ (default 0.01\,\AA).
  Sorted by mean convergence rate.}
\label{tab:v11}
\begin{tabular}{lcccc}
\toprule
Configuration & 0.5\,\AA & 1.0\,\AA & 1.5\,\AA & 2.0\,\AA \\
\midrule
\multicolumn{5}{l}{\emph{20k-step configs}} \\
RFO+Schl+PLS (20k) & \textbf{272} (94.8\%) & \textbf{245} (85.4\%) & \textbf{222} (77.4\%) & \textbf{183} (63.8\%) \\
GQT+Schl+PLS (20k) & 257 (89.5\%) & 220 (76.7\%) & 195 (67.9\%) & 167 (58.2\%) \\
\midrule
\multicolumn{5}{l}{\emph{10k-step configs (standard)}} \\
RFO+Schl+PLS        & 256 (89.2\%) & 225 (78.4\%) & 175 (61.0\%) & 124 (43.2\%) \\
RFO+PLS              & 255 (88.9\%) & 230 (80.1\%) & 176 (61.3\%) & 127 (44.3\%) \\
GQT+PLS              & 242 (84.3\%) & 191 (66.6\%) & 156 (54.4\%) & 118 (41.1\%) \\
GQT+Schl+PLS         & 239 (83.3\%) & 195 (67.9\%) & 154 (53.7\%) & 120 (41.8\%) \\
\midrule
\multicolumn{5}{l}{\emph{Trust floor = 0.005\,\AA}} \\
GQT+Schl+PLS+trf     & 206 (71.8\%) & 157 (54.7\%) & 125 (43.6\%) & 70 (24.4\%) \\
RFO+Schl+PLS+trf     & 201 (70.0\%) & 157 (54.7\%) & 89 (31.0\%) & 60 (20.9\%) \\
\bottomrule
\end{tabular}
\end{table}

\subsubsection{Factor Analysis}
\label{sec:v11-factors}

\paragraph{20k steps: the second breakthrough.}
Doubling the step budget from 10k to 20k produces the largest gains seen in
any intervention except PLS itself:

\begin{table}[h]
\centering
\caption{Effect of 20k steps vs.\ 10k (both with Schlegel+PLS).
  $\Delta$ = convergence gain in percentage points.}
\label{tab:20k-effect}
\begin{tabular}{lcccc}
\toprule
Optimizer & 0.5\,\AA & 1.0\,\AA & 1.5\,\AA & 2.0\,\AA \\
\midrule
RFO+Schl+PLS $\Delta$ & $+$5.6 & $+$7.0 & $+$16.4 & $+$20.6 \\
GQT+Schl+PLS $\Delta$ & $+$6.3 & $+$8.7 & $+$14.3 & $+$16.4 \\
\bottomrule
\end{tabular}
\end{table}

The effect grows dramatically with noise: $+$6\,pts at 0.5\,\AA\ but
$+$17--21\,pts at 2.0\,\AA.  At 2.0\,\AA, RFO+Schl+PLS+20k reaches
183/287 (63.8\%), a \textbf{+19.2\,pt improvement over any previous 10k
configuration} (v10c GQT+PLS: 128/287, 44.6\%).

\textbf{Which failure class benefits?}  Autopsy comparison for
RFO+Schl+PLS at each noise level (10k $\to$ 20k):

\begin{table}[h]
\centering
\caption{Failure class changes from 10k to 20k steps (RFO+Schlegel+PLS).
  Negative values indicate rescued samples.}
\label{tab:20k-autopsy}
\begin{tabular}{lrrrr}
\toprule
Failure class & $\Delta$(0.5\,\AA) & $\Delta$(1.0\,\AA) & $\Delta$(1.5\,\AA) & $\Delta$(2.0\,\AA) \\
\midrule
oscillating        & $-$13 & $-$9  & $-$37 & $-$46 \\
almost\_converged   & $\phantom{-}$0 & $-$9  & $-$7  & $-$3  \\
slow\_convergence   & $\phantom{-}$0 & $-$1  & $-$1  & $-$10 \\
ghost\_modes        & $-$3  & $-$1  & $-$2  & $\phantom{-}$0  \\
\midrule
Total rescued       & $+$16 & $+$20 & $+$47 & $+$59  \\
\bottomrule
\end{tabular}
\end{table}

Across all noise levels, \textbf{oscillating failures account for 105/142
(73.9\%) of rescued samples}.  Almost\_converged contributes 19/142 (13.4\%),
slow\_convergence 12/142 (8.5\%), and ghost\_modes 6/142 (4.2\%).
The 20k improvement is \textbf{primarily an oscillation-resolution effect}:
period-$\sim$8 oscillators that do not damp within 10k steps eventually
settle given sufficient time.  The blind-mode population
(almost\_converged + ghost\_modes at the trust radius floor) is largely
unaffected---these represent a qualitatively different failure mechanism
that more steps alone cannot resolve.

\textbf{Cost:} At $n=2.0$, mean steps increase from 4229 to 7526 ($1.78\times$),
wall time from 171\,s to 282\,s ($1.65\times$).  The 20.6\,pt convergence
gain at 2.0\,\AA\ justifies this cost for applications where convergence
rate matters more than throughput.

\paragraph{Step count distributions.}
The distribution of converged\_step reveals the convergence dynamics:

\begin{table}[h]
\centering
\caption{Step count statistics for converged samples (RFO+Schlegel+PLS).
  P90 = 90th percentile.  The median is much lower than the mean because
  most samples converge quickly; the mean is inflated by a long tail of
  slow oscillators.}
\label{tab:step-distributions}
\begin{tabular}{ccrrrr}
\toprule
Config & Noise & $N_\text{conv}$ & Mean & Median & P90 \\
\midrule
10k & 0.5\,\AA & 256 & 1{,}533 & 514 & 4{,}525 \\
    & 1.0\,\AA & 225 & 2{,}522 & 1{,}683 & 6{,}673 \\
    & 1.5\,\AA & 175 & 3{,}760 & 3{,}725 & 7{,}677 \\
    & 2.0\,\AA & 124 & 4{,}229 & 3{,}992 & 8{,}521 \\
\midrule
20k & 0.5\,\AA & 272 & 2{,}250 & 623 & 6{,}657 \\
    & 1.0\,\AA & 245 & 3{,}428 & 1{,}853 & 8{,}526 \\
    & 1.5\,\AA & 222 & 5{,}972 & 4{,}834 & 14{,}723 \\
    & 2.0\,\AA & 183 & 7{,}526 & 6{,}078 & 17{,}211 \\
\bottomrule
\end{tabular}
\end{table}

Of the 20k-converged samples, the fraction converging \emph{after} step
10{,}000 grows with noise: 5.9\% at $n=0.5$, 8.2\% at $n=1.0$, 21.2\%
at $n=1.5$, and 32.2\% at $n=2.0$.  At $n=2.0$, 59 samples converge
between steps 10k and 20k, with median convergence step 13{,}677 and maximum
19{,}908.  The converged-step distribution for these late convergers is roughly
uniform over the 10k--20k interval, \emph{not} concentrated near 10k.
This suggests many would benefit from even longer runs.

\textbf{Failure wall}: Nearly all non-converging samples in the 20k configs
hit the 20k step limit: 104/104 at $n=2.0$, 64/65 at $n=1.5$, 42/42 at
$n=1.0$.  No samples stopped early due to internal stagnation detection,
confirming these trajectories are actively evolving, not stuck.

\textbf{Mean step scaling}: Across the four standard 10k algorithms, mean
converged step scales approximately linearly with noise: $\sim$1{,}500 at
0.5\,\AA, $\sim$2{,}500 at 1.0\,\AA, $\sim$3{,}500 at 1.5\,\AA,
$\sim$4{,}000 at 2.0\,\AA.  RFO and GQT have similar mean steps at matched
noise levels ($\pm$10\%).

\paragraph{Trust radius floor 0.005\,\AA: catastrophic.}
Lowering the trust radius floor from 0.01\,\AA\ to 0.005\,\AA\ produces
\textbf{uniformly severe degradation}:

\begin{table}[h]
\centering
\caption{Effect of trust radius floor 0.005\,\AA\ vs.\ default 0.01\,\AA\
  (both with Schlegel+PLS, 10k steps).  All values are convergence change
  in percentage points.}
\label{tab:trf-effect}
\begin{tabular}{lcccc}
\toprule
Optimizer & 0.5\,\AA & 1.0\,\AA & 1.5\,\AA & 2.0\,\AA \\
\midrule
RFO+Schl+PLS $\Delta$ & $-$19.2 & $-$23.7 & $-$30.0 & $-$22.3 \\
GQT+Schl+PLS $\Delta$ & $-$11.5 & $-$13.2 & $-$10.1 & $-$17.4 \\
\bottomrule
\end{tabular}
\end{table}

RFO is hit harder than GQT ($-$19 to $-$30\,pts vs.\ $-$10 to $-$17\,pts).
Autopsy at $n=2.0$ for RFO+Schl+PLS (floor~0.01 $\to$ 0.005):
oscillating $117 \to 144$ ($+$27), slow\_convergence $15 \to 45$ ($+$30),
almost\_converged $14 \to 21$ ($+$7).

\textbf{Mechanism}: At floor~0.005\,\AA, each step displaces atoms half as
far.  This (a)~doubles the time needed to escape oscillation,
(b)~makes recovery from trust-radius collapse harder (requires twice as many
consecutive good steps to grow the radius back), and (c)~produces steps too
small to efficiently flip negative-eigenvalue modes.  The
slow\_convergence count \emph{tripling} ($15 \to 45$) is the clearest
signature: samples are making progress but at a rate insufficient to
converge within 10k steps.

\textbf{Conclusion}: The trust radius floor of 0.01\,\AA\ is already
optimal.  This value matches the $\varepsilon^{1/2}$ prediction
(Section~\ref{sec:eps-half}): with $\varepsilon = 10^{-4}$\,eV/\AA,
the floor should not be smaller than $\sqrt{\varepsilon} \approx 0.01$.

\paragraph{Schlegel TR: neutral when PLS is active.}
Comparing configs that differ only in Schlegel on/off (all with PLS):

\begin{table}[h]
\centering
\caption{Effect of Schlegel trust radius update when PLS is active.
  $\Delta$ = (with Schlegel) $-$ (without Schlegel), in percentage points.}
\label{tab:schlegel-with-pls}
\begin{tabular}{lcccc}
\toprule
Comparison & 0.5\,\AA & 1.0\,\AA & 1.5\,\AA & 2.0\,\AA \\
\midrule
RFO+Schl+PLS vs.\ RFO+PLS & $+$0.3 & $-$1.7 & $-$0.3 & $-$1.0 \\
GQT+Schl+PLS vs.\ GQT+PLS & $-$1.0 & $+$1.4 & $-$0.7 & $+$0.7 \\
\bottomrule
\end{tabular}
\end{table}

All deltas are within the $\pm$1--4\,pt seed variance
(Table~\ref{tab:reproducibility}).  \textbf{Schlegel TR has no detectable
effect when PLS is active.}  This is consistent with the interpretation
that PLS already handles step-length refinement, making the Schlegel
corrections to trust-radius growth/shrink rules redundant.

In v10c, Schlegel appeared to help because it was only tested in
combination with PLS (RFO+Schlegel+PLS was champion).  v11 reveals that
removing Schlegel from RFO+PLS produces \emph{identical or slightly better}
convergence, particularly at $n=1.0$ where RFO+PLS (230/287, 80.1\%)
outperforms RFO+Schlegel+PLS (225/287, 78.4\%) by 5 samples.

\paragraph{Interaction effects.}
Two-way interactions between the three factors (optimizer, 20k steps,
trust floor), computed as $\Delta(\text{A+B}) - \Delta(\text{A}) -
\Delta(\text{B})$:

\begin{table}[h]
\centering
\caption{Interaction terms (percentage points).  Positive = synergistic,
  negative = antagonistic.  All interactions measured with Schlegel+PLS fixed.}
\label{tab:interactions}
\begin{tabular}{lccccc}
\toprule
Interaction & 0.5\,\AA & 1.0\,\AA & 1.5\,\AA & 2.0\,\AA & Interpretation \\
\midrule
20k $\times$ optimizer  & $-$0.7 & $-$1.7 & $+$2.1 & $+$4.2 & 20k helps RFO slightly more at high noise \\
trf $\times$ optimizer  & $-$7.7 & $-$10.5 & $-$19.9 & $-$4.9 & trf hurts RFO much more than GQT \\
Schl $\times$ optimizer & $+$1.4 & $-$3.1 & $+$0.3 & $-$1.7 & No consistent interaction \\
\bottomrule
\end{tabular}
\end{table}

The trf$\times$optimizer interaction is the strongest: lowering the trust
floor hurts RFO 8--20\,pts more than GQT.  This is consistent with RFO's
augmented-Hessian shift producing larger step \emph{directions} that the
trust radius must clip; at a lower floor, the clipping is more aggressive,
disproportionately penalizing RFO.  The 20k$\times$optimizer interaction
is weak ($\pm$4\,pts), confirming that the 20k benefit is optimizer-independent.
The Schlegel$\times$optimizer interaction is negligible, as expected given
Schlegel's overall neutrality.

\paragraph{GQT+Schlegel+PLS: the never-tested combo.}
Adding Schlegel to GQT+PLS yields:
\begin{center}
\begin{tabular}{lcccc}
\toprule
Config & 0.5\,\AA & 1.0\,\AA & 1.5\,\AA & 2.0\,\AA \\
\midrule
GQT+PLS       & 242 (84.3\%) & 191 (66.6\%) & 156 (54.4\%) & 118 (41.1\%) \\
GQT+Schl+PLS  & 239 (83.3\%) & 195 (67.9\%) & 154 (53.7\%) & 120 (41.8\%) \\
$\Delta$       & $-$1.0       & $+$1.4       & $-$0.7       & $+$0.7       \\
\bottomrule
\end{tabular}
\end{center}
The effect is within noise.  GQT+Schlegel+PLS does \emph{not} close the gap
to RFO+PLS (which leads by 4.5--13.6\,pts at all noise levels).

\paragraph{RFO vs.\ GQT with PLS.}
Both optimizers tested with PLS (no Schlegel) for a clean comparison:
\begin{center}
\begin{tabular}{lcccc}
\toprule
Config & 0.5\,\AA & 1.0\,\AA & 1.5\,\AA & 2.0\,\AA \\
\midrule
RFO+PLS  & 255 (88.9\%) & 230 (80.1\%) & 176 (61.3\%) & 127 (44.3\%) \\
GQT+PLS  & 242 (84.3\%) & 191 (66.6\%) & 156 (54.4\%) & 118 (41.1\%) \\
$\Delta$ (RFO$-$GQT) & $+$4.5 & $+$13.6 & $+$7.0 & $+$3.1 \\
\bottomrule
\end{tabular}
\end{center}
\textbf{RFO wins at all noise levels}, with the largest advantage at
$n=1.0$ ($+$13.6\,pts).  The v10c crossover (GQT winning at 2.0\,\AA)
does not appear in this seed: $+$3.1\,pts favoring RFO.  Given the
$\pm$3.5\,pt seed variance (Table~\ref{tab:reproducibility}), the two
methods are statistically indistinguishable at 2.0\,\AA, but RFO
\emph{consistently} wins at 0.5--1.5\,\AA\ across both seeds.

The RFO advantage at $n=1.0$ is particularly striking: 39 additional
converged samples.  This is the noise level where both the Hessian
spectrum is sufficiently ill-conditioned to challenge the optimizer
($\kappa \sim 10^{4.4}$) but not so ill-conditioned that both methods
are equally lost ($\kappa \sim 10^{6.4}$ at 2.0\,\AA).

Mean steps at $n=2.0$: RFO+PLS = 4023 vs.\ GQT+PLS = 3873 (RFO takes
$\sim$4\% more steps but converges 3 more samples).

\subsubsection{v11 Failure Autopsy}
\label{sec:v11-autopsy}

\begin{table}[h]
\centering
\caption{v11 overall failure classification (9164 trajectories, all PLS, no ARC).}
\label{tab:v11-autopsy}
\begin{tabular}{lrr}
\toprule
Classification & Count & Fraction \\
\midrule
converged          & 5,800 & 63.3\% \\
oscillating        & 2,026 & 22.1\% \\
almost\_converged   &   917 & 10.0\% \\
ghost\_modes        &   306 &  3.3\% \\
slow\_convergence   &   111 &  1.2\% \\
drifting           &     4 &  0.0\% \\
\bottomrule
\end{tabular}
\end{table}

Compared to the v10c non-ARC failure landscape (oscillating $\sim$50\%,
almost\_converged $\sim$30\%, ghost\_modes $\sim$10\% of failures), the v11
distribution is inflated by the trust-floor configs (which generate
disproportionate oscillating and slow\_convergence failures).  Among the
six standard 10k configs (no trf, no 20k), the failure distribution
matches v10c within noise.

\textbf{Key autopsy insight}: The 20k configs shift the failure
distribution toward almost\_converged and ghost\_modes (the blind-mode
population), because oscillating failures---the largest class in 10k
configs---are resolved by additional steps.  At $n=2.0$ with 20k,
oscillating drops from 117 to 71 ($-$39\%), while ghost\_modes remains
fixed at 17 and almost\_converged drops only from 14 to 11.  This confirms
that \textbf{the residual failure population after 20k steps is dominated
by blind modes}, not oscillation.

\paragraph{Per-configuration failure breakdown.}
Grouping the 32 v11 configs by algorithm family, the per-noise failure
counts (Table~\ref{tab:per-config-autopsy}) show clear patterns:
(1)~trf configs have 2--3$\times$ more oscillating failures than their
default-floor counterparts; (2)~20k configs have 30--40\% fewer oscillating
failures; (3)~GQT has relatively more almost\_converged failures while
RFO has relatively more oscillating failures at matched noise/features.

\begin{table}[h]
\centering
\caption{Per-config autopsy summary for v11.  Columns: Conv = converged,
  Osc = oscillating, AC = almost\_converged, GM = ghost\_modes,
  SC = slow\_convergence.  Drifting ($\leq$4 total) omitted.
  $N \approx 286$ per config (varies by $\pm$2 due to SCINE errors).}
\label{tab:per-config-autopsy}
\begin{tabular}{llrrrrr}
\toprule
Algorithm & Noise & Conv & Osc & AC & GM & SC \\
\midrule
\multicolumn{7}{l}{\emph{RFO+Schlegel+PLS (20k)}} \\
 & 0.5 & 272 & 5 & 9 & 0 & 0 \\
 & 1.0 & 245 & 26 & 12 & 4 & 0 \\
 & 1.5 & 222 & 39 & 16 & 7 & 2 \\
 & 2.0 & 183 & 71 & 11 & 17 & 5 \\
\midrule
\multicolumn{7}{l}{\emph{GQT+Schlegel+PLS (20k)}} \\
 & 0.5 & 257 & 4 & 21 & 3 & 0 \\
 & 1.0 & 220 & 28 & 33 & 6 & 0 \\
 & 1.5 & 195 & 44 & 31 & 16 & 0 \\
 & 2.0 & 167 & 79 & 31 & 10 & 0 \\
\midrule
\multicolumn{7}{l}{\emph{RFO+Schlegel+PLS (10k)}} \\
 & 0.5 & 256 & 18 & 9 & 3 & 0 \\
 & 1.0 & 225 & 35 & 21 & 5 & 1 \\
 & 1.5 & 175 & 76 & 23 & 9 & 3 \\
 & 2.0 & 124 & 117 & 14 & 17 & 15 \\
\midrule
\multicolumn{7}{l}{\emph{RFO+PLS (10k)}} \\
 & 0.5 & 255 & 15 & 15 & 1 & 0 \\
 & 1.0 & 230 & 33 & 16 & 6 & 2 \\
 & 1.5 & 176 & 81 & 18 & 7 & 4 \\
 & 2.0 & 127 & 109 & 24 & 17 & 10 \\
\midrule
\multicolumn{7}{l}{\emph{GQT+PLS (10k)}} \\
 & 0.5 & 242 & 11 & 25 & 7 & 0 \\
 & 1.0 & 191 & 43 & 46 & 7 & 0 \\
 & 1.5 & 156 & 75 & 45 & 10 & 0 \\
 & 2.0 & 118 & 116 & 26 & 24 & 3 \\
\midrule
\multicolumn{7}{l}{\emph{GQT+Schlegel+PLS (10k)}} \\
 & 0.5 & 239 & 16 & 27 & 3 & 0 \\
 & 1.0 & 195 & 38 & 43 & 11 & 0 \\
 & 1.5 & 154 & 78 & 41 & 13 & 0 \\
 & 2.0 & 120 & 112 & 32 & 22 & 2 \\
\midrule
\multicolumn{7}{l}{\emph{RFO+Schlegel+PLS+trf (10k)}} \\
 & 0.5 & 201 & 53 & 26 & 5 & 1 \\
 & 1.0 & 157 & 85 & 36 & 7 & 2 \\
 & 1.5 & 89 & 125 & 47 & 11 & 14 \\
 & 2.0 & 60 & 144 & 21 & 17 & 45 \\
\midrule
\multicolumn{7}{l}{\emph{GQT+Schlegel+PLS+trf (10k)}} \\
 & 0.5 & 206 & 31 & 43 & 5 & 0 \\
 & 1.0 & 157 & 70 & 54 & 6 & 0 \\
 & 1.5 & 125 & 84 & 64 & 13 & 0 \\
 & 2.0 & 70 & 165 & 37 & 17 & 0 \\
\bottomrule
\end{tabular}
\end{table}

\paragraph{SCINE error analysis.}
Three samples crash with \texttt{UnboundLocalError: cannot access local
variable `out\_new'} --- a SCINE/Sparrow internal error during the PLS
step evaluation:
\begin{itemize}[nosep]
  \item Sample~16: crashes in all 8 configs at $n=0.5$ (both GQT and RFO).
    Converges at $n \geq 1.0$ in RFO configs.
  \item Sample~173: crashes in 4 GQT configs at $n=0.5$ only.
  \item Sample~236: crashes in all 8 configs at $n=1.5$ only.
\end{itemize}
The crash is deterministic, noise-level-specific, and affects the same
sample regardless of optimizer or step budget.  It appears to be triggered
by a specific molecular geometry encountered during the PLS evaluation
at those noise levels.  Total impact: 16 missing trajectories out of
9{,}184 ($<0.2\%$), negligible for aggregate statistics.

\subsubsection{Hardest Samples Across Runs}
\label{sec:hardest}

\begin{table}[h]
\centering
\caption{Top 10 hardest samples in v11 (converging in $\leq$10/32 configs).
  $n_\text{neg}^*$ and $\lambda_\text{min}^*$ are from the best failing attempt
  (highest $\lambda_\text{min}$ among failures).  v10c rate is out of 56 configs
  (includes non-PLS and ARC).  All share the blind-mode spectral signature:
  $|\lambda| < 10^{-3}$, gradient overlap $<0.1$.}
\label{tab:hardest-10}
\begin{tabular}{clccccc}
\toprule
Sample & Formula & v11 & v10c & $n_\text{neg}^*$ & $\lambda_\text{min}^*$ & Overlap$^*$ \\
\midrule
179 & C$_3$H$_8$N$_2$O & 5/32 (16\%) & 8/56 (14\%) & 3 & $-8.6\times10^{-5}$ & 0.049 \\
199 & C$_5$H$_5$NO      & 6/32 (19\%) & 16/56 (29\%) & 5 & $-8\times10^{-6}$ & 0.022 \\
129 & C$_3$H$_5$NO$_2$  & 7/32 (22\%) & 20/56 (36\%) & 2 & $-3.4\times10^{-5}$ & 0.068 \\
16  & C$_3$H$_5$NO$_2$  & 9/32 (28\%) & 7/56 (13\%) & 1 & $-5.3\times10^{-5}$ & 0.013 \\
100 & C$_3$H$_5$NO$_2$  & 9/32 (28\%) & 12/56 (21\%) & 1 & $-4\times10^{-6}$ & 0.023 \\
147 & C$_3$H$_5$NO$_2$  & 9/32 (28\%) & 12/56 (21\%) & 1 & $-2\times10^{-6}$ & 0.079 \\
183 & C$_3$H$_8$N$_2$O  & 9/32 (28\%) & 6/56 (11\%) & 2 & $-7\times10^{-6}$ & 0.015 \\
193 & C$_5$H$_5$NO      & 10/32 (31\%) & 20/56 (36\%) & 1 & $-1.2\times10^{-4}$ & 0.033 \\
228 & C$_5$H$_5$NO      & 10/32 (31\%) & 10/56 (18\%) & 2 & $-1.2\times10^{-4}$ & 0.003 \\
277 & C$_5$H$_5$NO      & 10/32 (31\%) & 7/56 (13\%) & 2 & $-1\times10^{-6}$ & 0.040 \\
\bottomrule
\end{tabular}
\end{table}

\textbf{Molecular pattern}: Only three molecular formulae appear:
C$_3$H$_5$NO$_2$ (samples 16, 100, 129, 147), C$_5$H$_5$NO (199, 193,
228, 277), and C$_3$H$_8$N$_2$O (179, 183).  These are all nitrogen- and
oxygen-containing small organic molecules, suggesting the Hessian noise
floor is particularly problematic near N--O bond rearrangement coordinates.

Sample~179 is hardest in both v10c (14\%) and v11 (16\%).  Sample~16
worsens from v10c (13\%) relative to the expanded v11 grid because its
SCINE crash at $n=0.5$ (8 of 8 GQT configs) effectively removes half
the grid.  Samples~184 and~210, the hardest in v10c (4/56, 7\% each),
are \emph{not} among v11's hardest, confirming they benefit from the all-PLS grid.

All 10 samples converge only at low noise levels: 9/10 have their
best convergence at $n=0.5$; only sample~16 converges more readily
at $n=1.0$--$2.0$ (9 configs) than at $n=0.5$ (0 configs, due to SCINE
crashes in GQT and the sample being intrinsically hard for GQT).

\subsubsection{Rescued by 20k Steps vs.\ Permanently Hard}
\label{sec:rescued-vs-hard}

For the champion configuration (RFO+Schlegel+PLS), every sample at each
noise level was classified as \emph{easy} (converges at 10k), \emph{rescued
by 20k} (fails at 10k, converges at 20k), or \emph{permanently hard}
(fails at both 10k and 20k):

\begin{table}[h]
\centering
\caption{Sample classification by 10k/20k convergence (RFO+Schlegel+PLS).
  $N=287$ samples per noise level (minus SCINE errors).}
\label{tab:rescued-vs-hard}
\begin{tabular}{ccccc}
\toprule
Noise & Easy (10k) & Rescued (20k) & Perm.\ hard & Rescued median step \\
\midrule
0.5\,\AA & 256 & 16 & 15 & 13{,}280 \\
1.0\,\AA & 225 & 20 & 42 & 13{,}312 \\
1.5\,\AA & 175 & 47 & 65 & 14{,}604 \\
2.0\,\AA & 124 & 59 & 104 & 13{,}677 \\
\bottomrule
\end{tabular}
\end{table}

The rescued population grows from 16 at $n=0.5$ to 59 at $n=2.0$, while the
permanently hard population grows from 15 to 104.  Median convergence step
for rescued samples is remarkably constant at $\sim$13{,}500 across all noise
levels, suggesting a common oscillation-damping timescale.

At $n=2.0$, the 59 rescued samples span a wide convergence range:
6 converge just past 10k (steps 10{,}021--10{,}538), while 5 converge near
the 20k wall (steps 19{,}260--19{,}908), including sample~100 at step~19{,}908.
This implies that even 20k steps barely suffices for the slowest oscillators.

\textbf{Verification}: 10k and 20k trajectories are identical for the first
10k steps (confirmed by energy comparison at steps 0, 100, 1000, 5000, 9999
for samples 179 and 183).  The 20k benefit is purely from extra steps, not
from a different trajectory.

\textbf{Failure wall}: Nearly all permanently hard samples hit the 20k step
limit (104/104 at $n=2.0$, 64/65 at $n=1.5$).  None stop early due to
stagnation detection---they are actively making progress, just too slowly
to converge.

\subsubsection{Trajectory Analysis}
\label{sec:trajectory-analysis}

Diagnostic trajectory data (per-step eigenvalues, forces, trust radius, PLS
activity) was extracted for key samples at $n=2.0$ with RFO+Schlegel+PLS.
Three representative trajectories illustrate the failure mechanisms:

\paragraph{Sample 179 (permanently hard, 5/32 converges).}
At $n=2.0$: starts with $\lambda_\text{min}=-8.25$, 16 negative eigenvalues,
force = 1.59\,eV/\AA.  After 10k steps: $\lambda_\text{min}=-0.090$,
9 negative eigenvalues, force = 0.036.  After 20k: $\lambda_\text{min}=-0.091$,
8 negative eigenvalues, force = 0.034.  \textbf{No eigenvalue zero-crossings}
in either 10k or 20k trajectory.  Trust radius at floor (0.01\,\AA) for
97--98\% of the trajectory.  PLS applied in 730/10{,}000 (7.3\%) and
1{,}716/20{,}000 (8.6\%) steps.  This is not oscillation but \textbf{genuine
stagnation in a deep saddle basin}: the molecule has strongly negative
eigenvalues that the optimizer slowly reduces from $-8.25$ to $-0.09$ over
20k steps, but never reaches zero.  Even 100k steps would likely be
insufficient without a qualitatively different escape mechanism.

\paragraph{Sample 183 (rescued at step 17{,}560).}
At $n=2.0$: starts with $\lambda_\text{min}=-3.88$, 13 negative eigenvalues,
force = 1.78.  At step 10{,}000 (end of 10k run): $\lambda_\text{min}=-4.85
\times10^{-4}$, 4 negative eigenvalues, force = 0.018.  With 20k: converges
at step 17{,}560 when $\lambda_\text{min}$ crosses zero.  PLS is heavily
active: applied in 6{,}053/17{,}561 (34.5\%) of steps.  Trust radius at floor
for 90\% of the trajectory.  This sample represents the \textbf{slow
oscillation-damping} profile: negative eigenvalues slowly decrease in
magnitude from $O(1)$ to $O(10^{-4})$ over $\sim$10k steps, then eventually
cross zero by step $\sim$17k.

\paragraph{Sample 228 (rescued at step 16{,}563).}
Similar to 183: starts at $\lambda_\text{min}=-8.54$, 12 negative eigenvalues.
At step 10{,}000: $\lambda_\text{min}=-1.3\times10^{-3}$, 3 negative
eigenvalues.  Converges at step 16{,}563 when eigenvalues finally cross zero.
PLS applied 24\% of steps.

\textbf{Common pattern}: The rescued samples share a profile of monotonic
eigenvalue improvement (not oscillation in the classical sense) from
$O(1)$ negative eigenvalues to $O(10^{-4})$ over $\sim$10k steps, then
slow final convergence to zero over an additional 3--10k steps.  The
permanently hard samples (sample~179) have this same trajectory but at a
slower rate, never reaching the zero crossing within 20k steps.

\subsubsection{Consolidated Best-Ever Table}
\label{sec:best-ever}

\begin{table}[h]
\centering
\caption{Best strict convergence rate at each noise level across all runs
  and seeds (287 samples).  10k and 20k results separated.}
\label{tab:best-ever}
\begin{tabular}{ccll}
\toprule
Noise & Best rate & Configuration & Run \\
\midrule
\multicolumn{4}{l}{\emph{10k steps}} \\
0.5\,\AA & 256/287 (89.2\%) & RFO+Schl+PLS       & v11 \\
1.0\,\AA & 230/287 (80.1\%) & RFO+PLS             & v11 \\
1.5\,\AA & 182/287 (63.4\%) & RFO+Schl+PLS        & v10c \\
2.0\,\AA & 128/287 (44.6\%) & GQT+PLS             & v10c \\
\midrule
\multicolumn{4}{l}{\emph{20k steps}} \\
0.5\,\AA & 272/287 (94.8\%) & RFO+Schl+PLS (20k)  & v11 \\
1.0\,\AA & 245/287 (85.4\%) & RFO+Schl+PLS (20k)  & v11 \\
1.5\,\AA & 222/287 (77.4\%) & RFO+Schl+PLS (20k)  & v11 \\
2.0\,\AA & 183/287 (63.8\%) & RFO+Schl+PLS (20k)  & v11 \\
       & \multicolumn{3}{l}{\emph{Superseded: see Table~\ref{tab:best-ever-v12b} for v12/v12b best-ever}} \\
\bottomrule
\end{tabular}
\end{table}

Note: v10c and v11 use different random seeds.  The v10c entries at 1.5 and
2.0\,\AA\ may reflect seed-favorable outcomes ($\pm$3.5\,pts); within-seed
v11 10k results at those noise levels are 175/287 (61.0\%) and 127/287
(44.3\%).

%=============================================================================
\subsection{v12 Results (mar8 Run, March 2026)}
\label{sec:v12}
%=============================================================================

The v12 grid tested two orthogonal interventions against the remaining
failure population identified in v11: (1)~perturbation kicks targeting
oscillating and blind-mode failures, and (2)~extending the step budget
to 50k.  All configs use RFO+PLS (the v11 recommended baseline) with
trust floor 0.01\,\AA.  A separate ``v12\_quick'' run tested kick variants
at $n=2.0$\,\AA.

\subsubsection{v12 Convergence Table}
\label{sec:v12-table}

\begin{table}[h]
\centering
\caption{v12 convergence results ($N=287$, strict $n_\text{neg}=0$, all
  configs use RFO+PLS).  ``osc'' = oscillation kick enabled,
  ``blind'' = blind-mode kick enabled, ``e1'' = kick along 2nd-lowest
  eigenvector instead of lowest.}
\label{tab:v12}
\begin{tabular}{lcccc}
\toprule
Configuration & 0.5\,\AA & 1.0\,\AA & 1.5\,\AA & 2.0\,\AA \\
\midrule
\multicolumn{5}{l}{\emph{50k-step configs}} \\
RFO+PLS (50k)          & 273 (95.1\%) & \textbf{275} (95.8\%) & \textbf{266} (92.7\%) & --- \\
\midrule
\multicolumn{5}{l}{\emph{20k-step configs}} \\
RFO+PLS (20k, baseline) & \textbf{267} (93.0\%) & 256 (89.2\%) & 230 (80.1\%) & 175 (61.0\%) \\
RFO+PLS+osc\_kick       & 266 (92.7\%) & 258 (89.9\%) & ---          & --- \\
RFO+PLS+blind\_kick     & 266 (92.7\%) & 256 (89.2\%) & ---          & --- \\
RFO+PLS+both\_kicks     & 266 (92.7\%) & 256 (89.2\%) & ---          & \textbf{192} (66.9\%) \\
RFO+PLS+both\_kicks+e1  & 265 (92.3\%) & 256 (89.2\%) & ---          & 191 (66.6\%) \\
\bottomrule
\end{tabular}
\end{table}

Nine samples at $n=0.5$ and two at $n=1.0$ produced SCINE errors
(\texttt{out\_new} variable bug), consistent across all kick variants
at each noise level.

\subsubsection{50k Steps: The Major Finding}
\label{sec:v12-50k}

Extending the step budget from 20k to 50k is a pure, lossless gain at
every tested noise level:

\begin{table}[h]
\centering
\caption{20k vs.\ 50k step comparison (RFO+PLS).  ``Rescued'' = samples
  that converge with 50k but not 20k.  ``Lost'' = samples that converge
  with 20k but not 50k.}
\label{tab:v12-50k}
\begin{tabular}{ccccccc}
\toprule
Noise & 20k conv & 50k conv & $\Delta$ (pts) & Rescued & Lost & Mean conv.\ step (rescued) \\
\midrule
0.5\,\AA & 267 & 273 & $+$2.1 &  6 & 0 & 32{,}977 \\
1.0\,\AA & 256 & 275 & $+$6.6 & 19 & 0 & 27{,}423 \\
1.5\,\AA & 230 & 266 & $+$12.5 & 36 & 0 & 28{,}041 \\
\bottomrule
\end{tabular}
\end{table}

All rescued samples required $>$20k steps (minimum 20{,}153; maximum
48{,}649).  No sample was lost, confirming 50k is a strict superset of 20k
for these configurations.  The gains grow with noise level: $+$2.1\,pts at
0.5\,\AA\ vs.\ $+$12.5\,pts at 1.5\,\AA, consistent with higher-noise
trajectories requiring more iterations to resolve oscillating modes.

The remaining failures at 50k are hard-core: 5 true failures $+$ 9
SCINE errors at $n=0.5$; 10 true failures $+$ 2 errors at $n=1.0$; 21
true failures at $n=1.5$.  These show persistent negative eigenvalues
(often $\lambda_\text{min} < -0.01$) and elevated force norms even after
50{,}000 steps, indicating they are structurally stuck rather than merely
slow.

\subsubsection{ModeTracker}
\label{sec:v12-modetracker}

The kick and late-escape mechanisms introduced in v12/v12b rely on a
\texttt{ModeTracker} class that tracks the identity of negative
eigenmodes across optimization steps via eigenvector overlap matching.
At each step, the tracker matches current negative eigenvectors to
previously tracked modes using a greedy assignment based on
$|\langle v_i^{(k)}, v_j^{(k-1)}\rangle|$, with a minimum overlap
threshold of~0.5 to avoid spurious matches.  Each tracked mode
maintains: the current eigenvector and eigenvalue, the \emph{age}
(number of consecutive steps the mode has remained negative), the
\emph{blind age} (consecutive steps with
$|\mathbf{g} \cdot \mathbf{v}| / \|\mathbf{g}\| < \text{threshold}$),
and the gradient overlap $|\mathbf{g} \cdot \mathbf{v}| / \|\mathbf{g}\|$.
The tracker exposes two key queries:
\texttt{longest\_stuck()} returns the mode with the largest age
(continuously negative for the most steps), and
\texttt{longest\_blind()} returns the mode with the largest blind age.
These are used by the kick mechanisms (Section~\ref{sec:v12-kicks}) and
late escape (Section~\ref{sec:v12b}) to select which mode to perturb
along.

\subsubsection{Kick Mechanism Analysis}
\label{sec:v12-kicks}

v12 introduced two perturbation mechanisms to escape failure modes:
\begin{itemize}[nosep]
  \item \textbf{Oscillation kick}: perturbs geometry along the lowest
    eigenvector when oscillation is detected (patience = 3 steps,
    cooldown = 50 steps, scale = 0.1).
  \item \textbf{Blind-mode kick}: perturbs geometry when a negative
    eigenvalue has low gradient overlap ($<0.1$), targeting ghost modes
    (patience = 100 steps, scale = 0.5).
\end{itemize}

\paragraph{Kick magnitudes are very small.}
The perturbation displacement is $\texttt{scale} \times \texttt{trust\_floor}$,
giving: oscillation kicks $= 0.1 \times 0.01 = 0.001$\,\AA, and
blind-mode kicks $= 0.5 \times 0.01 = 0.005$\,\AA.  Typical atomic
displacements needed to escape a saddle basin are of order
$0.01$--$0.1$\,\AA, so v12's kicks are $10$--$100\times$ smaller than
what is physically needed.  This explains why kicks fire prolifically
($\sim$122 oscillation kicks per trajectory) yet fail to reduce oscillating
failures: each individual kick is too small to move the geometry out of
the basin of attraction.

\paragraph{At $n \leq 1.5$\,\AA, kicks never fired.}
All kick counters (total\_osc\_kicks, total\_blind\_kicks) are zero
across all 287 samples at $n=0.5$, 1.0, and 1.5\,\AA.  The kick
thresholds were not met at these noise levels, so all kick variants
are identical to baseline (within $\pm$1 sample from stochastic effects).

\paragraph{At $n=2.0$\,\AA, kicks fired heavily and helped modestly.}
For the 13 failure trajectories with saved diagnostics, the both\_kicks
config averaged $\sim$122 oscillation kicks and $\sim$44 blind-mode kicks
per trajectory.  The net effect:

\begin{table}[h]
\centering
\caption{Failure class breakdown at $n=2.0$\,\AA\ (from v12\_quick autopsy).
  ``Osc'' = oscillating, ``AC'' = almost\_converged, ``GM'' = ghost\_modes,
  ``SC'' = slow\_convergence.}
\label{tab:v12-kicks-autopsy}
\begin{tabular}{lccccc}
\toprule
Configuration & Conv & Osc & AC & GM & SC \\
\midrule
RFO+PLS (baseline)  & 175 & 71 & 19 & 12 &  8 \\
RFO+PLS+both\_kicks & 192 & 69 & 16 &  9 &  1 \\
RFO+PLS+both+e1     & 191 & 67 & 16 &  8 &  4 \\
\bottomrule
\end{tabular}
\end{table}

The both\_kicks config rescued 32 samples vs.\ baseline but also lost 15,
for a net gain of $+$17 ($+$5.9\,pts).  Of the 32 rescued, 18 were
oscillating, 10 almost\_converged, 2 slow\_convergence, and 1 each
ghost\_modes and drifting.  Of the 15 lost, 9 became oscillating, 4
almost\_converged, and 2 ghost\_modes---indicating kicks can destabilize
borderline trajectories while rescuing others.

The eigvec1 variant (kicking along the \emph{second}-lowest eigenvector)
shows essentially identical results (191 vs.\ 192), with a similar
rescue/loss pattern (30 rescued, 14 lost).  This suggests the kick
direction is not critical---the perturbation's main value is breaking
symmetry rather than targeting a specific mode.

\paragraph{Oscillation not specifically reduced.}
Despite the osc\_kick mechanism, oscillating failures dropped only
from 71 to 69 (both\_kicks) or 67 (eigvec1).  The kicks rescue some
oscillating samples but create new oscillating failures in roughly equal
proportion.  The net oscillating count is essentially unchanged.

\subsubsection{Comparison to v11}
\label{sec:v12-vs-v11}

\begin{table}[h]
\centering
\caption{v12 vs.\ v11 best (RFO+Schlegel+PLS, 20k).  Note v12 uses a
  different seed than v11; differences at 20k are within seed variance.}
\label{tab:v12-vs-v11}
\begin{tabular}{ccccc}
\toprule
Noise & v11 best (20k) & v12 baseline (20k) & v12 best kick (20k) & v12 (50k) \\
\midrule
0.5\,\AA & 272 (94.8\%) & 267 (93.0\%) & --- & 273 (95.1\%) \\
1.0\,\AA & 245 (85.4\%) & 256 (89.2\%) & 258 (89.9\%) & \textbf{275 (95.8\%)} \\
1.5\,\AA & 222 (77.4\%) & 230 (80.1\%) & --- & \textbf{266 (92.7\%)} \\
2.0\,\AA & 183 (63.8\%) & 175 (61.0\%) & 192 (66.9\%) & --- \\
\bottomrule
\end{tabular}
\end{table}

The 50k results at $n=1.0$ and $n=1.5$ are new best-ever rates by large
margins: $+$10.4\,pts at 1.0\,\AA\ and $+$15.3\,pts at 1.5\,\AA\ over
v11's 20k results.  However, cross-seed comparison means the true
improvement from 20k$\to$50k is better measured within-seed from
Table~\ref{tab:v12-50k}: $+$6.6\,pts and $+$12.5\,pts respectively.

At $n=0.5$, the 50k rate (95.1\%) matches v11's 20k rate (94.8\%) within
seed variance, and 9 of the 14 remaining failures are SCINE errors, not
optimizer failures.  This noise level is effectively \textbf{solved}
(the 5 true failures at 50k may be intrinsically unconvergeable under
DFTB0).

\subsubsection{Updated Best-Ever Table}
\label{sec:v12-best-ever}

\begin{table}[h]
\centering
\caption{Best strict convergence rate at each noise level across all runs
  and seeds ($N=287$).  Updated through v12.}
\label{tab:best-ever-v12}
\begin{tabular}{ccll}
\toprule
Noise & Best rate & Configuration & Run \\
\midrule
\multicolumn{4}{l}{\emph{10k steps}} \\
0.5\,\AA & 256/287 (89.2\%) & RFO+Schl+PLS       & v11 \\
1.0\,\AA & 230/287 (80.1\%) & RFO+PLS             & v11 \\
1.5\,\AA & 182/287 (63.4\%) & RFO+Schl+PLS        & v10c \\
2.0\,\AA & 128/287 (44.6\%) & GQT+PLS             & v10c \\
\midrule
\multicolumn{4}{l}{\emph{20k steps}} \\
0.5\,\AA & 272/287 (94.8\%) & RFO+Schl+PLS (20k)  & v11 \\
1.0\,\AA & 256/287 (89.2\%) & RFO+PLS (20k)        & v12 \\
1.5\,\AA & 230/287 (80.1\%) & RFO+PLS (20k)        & v12 \\
2.0\,\AA & 192/287 (66.9\%) & RFO+PLS+both\_kicks (20k) & v12 \\
       & \multicolumn{3}{l}{\emph{Superseded by v12b: 218/287 (76.0\%), see Table~\ref{tab:best-ever-v12b}}} \\
\midrule
\multicolumn{4}{l}{\emph{50k steps}} \\
0.5\,\AA & 273/287 (95.1\%) & RFO+PLS (50k)  & v12 \\
1.0\,\AA & 275/287 (95.8\%) & RFO+PLS (50k)  & v12 \\
1.5\,\AA & 266/287 (92.7\%) & RFO+PLS (50k)  & v12 \\
\bottomrule
\end{tabular}
\end{table}

\subsubsection{v12 Summary and Next Steps}
\label{sec:v12-summary}

\textbf{Main conclusions:}
\begin{enumerate}[nosep]
  \item \textbf{50k steps is the dominant intervention}: $+$2 to $+$13\,pts
    over 20k, lossless (no samples lost), with all rescued samples needing
    $>$20k iterations.
  \item \textbf{Kicks are a no-op at $n \leq 1.5$\,\AA}: the perturbation
    thresholds were never met, leaving convergence rates identical to
    baseline.
  \item \textbf{Kicks give modest help at $n=2.0$\,\AA}: $+$5.9\,pts net,
    but high churn (32 rescued, 15 lost).  Kick direction (eigvec0 vs.\
    eigvec1) is irrelevant.
  \item \textbf{$n=0.5$\,\AA\ is effectively solved}: 95.1\% at 50k, with
    9/14 remaining failures being SCINE errors.
  \item \textbf{SCINE \texttt{out\_new} bug}: 9 errors at $n=0.5$ and 2 at
    $n=1.0$ from a calculator-level bug.  Fixing this could add $\sim$3\,pts
    at $n=0.5$.
\end{enumerate}

The v12 kick mechanism as implemented (fixed thresholds, fixed scale) is
too conservative at low noise (never triggers) and too aggressive at high
noise (destabilizes borderline samples).  Potential v12b improvements:
\begin{itemize}[nosep]
  \item \textbf{Adaptive kick scaling}: scale perturbation magnitude as
    $C \cdot |\lambda_i|^{1/2}$ rather than using a fixed constant.
  \item \textbf{Probe-based kicks}: probe a distance ladder in both $\pm$
    directions and accept the first displacement that reduces $n_\text{neg}$.
  \item \textbf{Late-onset kicks}: only activate after $>$10k steps to avoid
    destabilizing early-stage convergence.
  \item \textbf{50k + kicks at $n=2.0$}: combining the two interventions
    (not yet tested).
\end{itemize}

%=============================================================================
\subsection{v12b Results (mar8 Run, March 2026)}
\label{sec:v12b}
%=============================================================================

The v12b grid (``v12b\_quick'' run, job~1146854) tested three incremental
improvements to the v12 kick mechanism, all at $n=2.0$\,\AA\ with 20k
steps ($N=287$).  The v12 analysis (Section~\ref{sec:v12-summary})
identified that fixed-threshold, fixed-scale kicks are too conservative
at low noise and too aggressive at high noise.  v12b addresses this with:
\begin{enumerate}[nosep]
  \item \textbf{Adaptive kick scaling} (\texttt{adaptive\_kick\_scale}):
    scale perturbation magnitude as $C \cdot |\lambda_i|^{1/2}$,
    clamped to $[\texttt{trust\_floor},\; 10 \times \texttt{trust\_floor}]$
    (i.e., $[0.01, 0.1]$\,\AA), with default $C=0.1$.
    For example, at $|\lambda|=0.09$ (sample~179):
    $\text{mag} = 0.1 \times 0.3 = 0.03$\,\AA---$30\times$ larger than
    v12's fixed $0.001$\,\AA\ oscillation kick.
  \item \textbf{Probe-based kicks} (\texttt{blind\_kick\_probe}):
    instead of a single fixed perturbation, probes a distance ladder
    at $[1, 2, 5, 10, 20] \times \texttt{trust\_floor}$
    (i.e., $0.01$--$0.2$\,\AA), testing both $+$ and $-$ directions
    along the eigenvector at each distance.  The probe accepts the
    \emph{first} displacement where the number of negative eigenvalues
    ($n_\text{neg}$) decreases; lower energy is used only as a
    fallback if no $n_\text{neg}$ improvement is found at any
    probe point.  This is fundamentally different from
    ``test both directions and pick lower energy'': the primary
    acceptance criterion is spectral improvement, not energetic.
  \item \textbf{Late escape} (\texttt{late\_escape}):
    an additional perturbation mechanism activated only after 15{,}000
    steps, targeting trajectories that are stuck but not yet detected
    by the oscillation or blind-mode criteria ($\alpha=0.1$,
    cooldown = 500 steps).  Mode selection uses
    \texttt{longest\_stuck()} from the ModeTracker
    (Section~\ref{sec:v12-modetracker}): the mode that has been
    continuously negative for the most steps, rather than the
    instantaneously most-negative eigenvalue.  The displacement is
    accepted only if $n_\text{neg}$ decreases (conditional acceptance,
    not unconditional); if the perturbation does not reduce the number
    of negative eigenvalues, it is rejected and the geometry is
    unchanged.
\end{enumerate}

\subsubsection{v12b Convergence Table}
\label{sec:v12b-table}

\begin{table}[h]
\centering
\caption{v12b convergence results at $n=2.0$\,\AA\ ($N=287$, strict
  $n_\text{neg}=0$, all configs use RFO+PLS, 20k steps).  Features
  are cumulative: each row adds one feature to the previous.}
\label{tab:v12b}
\begin{tabular}{lccc}
\toprule
Configuration & Conv & Rate & $\Delta$ vs.\ baseline \\
\midrule
RFO+PLS (baseline)               & 180 & 62.7\% & ---        \\
+both\_kicks                      & 180 & 62.7\% & $+$0.0\,pts \\
+adaptive scaling                 & 181 & 63.1\% & $+$0.3\,pts \\
+probe kicks                      & 207 & 72.1\% & $+$9.4\,pts \\
+late escape (\textbf{full v12b}) & \textbf{218} & \textbf{76.0\%} & $+$13.2\,pts \\
\bottomrule
\end{tabular}
\end{table}

\subsubsection{Incremental Feature Analysis}
\label{sec:v12b-incremental}

Because features are cumulative, we can isolate each contribution by
comparing adjacent configs.  The rescue/loss analysis counts samples that
change convergence status between consecutive configs:

\begin{table}[h]
\centering
\caption{Incremental rescue/loss analysis.  Each row shows the effect of
  adding one feature to the previous config.  ``Rescued'' = newly converged,
  ``Lost'' = newly failed.}
\label{tab:v12b-incremental}
\begin{tabular}{lcccc}
\toprule
Feature added & Rescued & Lost & Net & $\Delta$ conv \\
\midrule
both\_kicks (vs.\ baseline) &  21 & 21 &   0 & $180 \to 180$ \\
adaptive scaling (vs.\ both\_kicks) & 27 & 26 & $+$1 & $180 \to 181$ \\
probe kicks (vs.\ adaptive) & 44 & 18 & $+$26 & $181 \to 207$ \\
late escape (vs.\ adapt+probe) & 16 &  5 & $+$11 & $207 \to 218$ \\
\bottomrule
\end{tabular}
\end{table}

\paragraph{Both\_kicks alone is a wash (net zero).}
This is a striking departure from v12's result at $n=2.0$, where
both\_kicks showed 192/287 vs.\ 175/287 baseline.  The difference is
the seed: v12b\_quick uses job~1146854 (different from v12\_quick's seed),
and the both\_kicks mechanism has high churn---21 rescued but 21 lost.
This confirms the v12 finding that fixed-scale kicks destabilize
borderline samples as often as they rescue stuck ones.

\paragraph{Adaptive scaling is marginal ($+$1 net).}
Scaling kick magnitude to the eigenvalue does not materially reduce
churn: 27 rescued vs.\ 26 lost.  The adaptive mechanism changes
\emph{which} samples are rescued/lost but not the overall balance.

\paragraph{Probe kicks are the breakthrough ($+$26 net).}
Probing a distance ladder ($0.01$--$0.2$\,\AA) in both $\pm$ directions
and accepting the first displacement that reduces $n_\text{neg}$
is the single most impactful feature in v12b.  It rescues 44 samples
while losing only 18---a 2.4:1 rescue-to-loss ratio, far better than the
$\sim$1:1 ratio of fixed kicks.  Of the 44 rescued samples (from the
\emph{adaptive} config's failure classes), 34 were oscillating, 5 were
almost\_converged, 3 were ghost\_modes, and 2 were slow\_convergence.
The probe mechanism primarily resolves oscillating failures by ensuring
the perturbation improves the spectral structure (fewer negative
eigenvalues), with lower energy as a fallback criterion.

\paragraph{Late escape adds a clean $+$11.}
The late-escape mechanism (activating after 15{,}000 steps) rescues 16
samples while losing only 5---a 3.2:1 ratio, the best of any feature.
The high rescue-to-loss ratio makes sense: by step 15{,}000, trajectories
that will converge naturally have already done so, so perturbations at
this stage are unlikely to destabilize converging samples.  All 5 losses
became oscillating, suggesting a small fraction of late-stage trajectories
are in a delicate equilibrium that the escape disrupts.

\subsubsection{Failure Autopsy Comparison}
\label{sec:v12b-autopsy}

\begin{table}[h]
\centering
\caption{Failure class breakdown at $n=2.0$\,\AA\ across v12b configs.
  ``Osc'' = oscillating, ``AC'' = almost\_converged, ``GM'' = ghost\_modes,
  ``SC'' = slow\_convergence, ``Dr'' = drifting.}
\label{tab:v12b-autopsy}
\begin{tabular}{lcccccc}
\toprule
Configuration & Conv & Osc & AC & GM & SC & Dr \\
\midrule
Baseline            & 180 & 68 & 17 & 13 &  9 & 0 \\
+both\_kicks         & 180 & 76 & 16 &  8 &  6 & 1 \\
+adaptive            & 181 & 80 & 13 &  7 &  6 & 0 \\
+probe               & 207 & 49 & 13 &  9 &  9 & 0 \\
+late (\textbf{v12b}) & \textbf{218} & 54 &  6 &  4 &  5 & 0 \\
\bottomrule
\end{tabular}
\end{table}

Several patterns emerge:

\begin{itemize}[nosep]
  \item \textbf{Oscillating failures decrease dramatically with probe}:
    $80 \to 49$ (adaptive$\to$probe), confirming probe kicks resolve
    oscillation.  Late escape slightly increases oscillating count
    ($49 \to 54$) because 5 previously non-oscillating trajectories
    become oscillating after the late perturbation, but net convergence
    still improves by~11.
  \item \textbf{Almost\_converged slashed by late escape}:
    $13 \to 6$.  Late-stage perturbation is exactly what
    almost\_converged samples need---they are near a minimum but
    stuck with residual negative eigenvalues.
  \item \textbf{Ghost modes halved overall}: $13 \to 4$ (baseline to
    full v12b).  The probe mechanism resolves some ghost modes by
    testing the sign of the perturbation direction.
  \item \textbf{69 total failures remain} in full v12b: 54 oscillating,
    6 almost\_converged, 4 ghost\_modes, 5 slow\_convergence.
    Oscillating remains the dominant failure mode (78\% of failures).
\end{itemize}

\subsubsection{Comparison to v12 and v11}
\label{sec:v12b-vs-prior}

\begin{table}[h]
\centering
\caption{Best convergence at $n=2.0$\,\AA\ across versions (20k steps,
  $N=287$).  Different runs use different seeds; within-seed comparisons
  are more reliable.}
\label{tab:v12b-vs-prior}
\begin{tabular}{llcc}
\toprule
Version & Best config at 2.0\,\AA & Conv & Rate \\
\midrule
v11     & RFO+Schl+PLS (20k)          & 183 & 63.8\% \\
v12     & RFO+PLS+both\_kicks (20k)   & 192 & 66.9\% \\
v12b    & RFO+PLS+adaptive+probe+late & \textbf{218} & \textbf{76.0\%} \\
\midrule
\multicolumn{4}{l}{\emph{Within-seed comparison (v12b run, job 1146854):}} \\
\midrule
v12b baseline & RFO+PLS (20k)           & 180 & 62.7\% \\
v12b best     & +adaptive+probe+late     & 218 & 76.0\% \\
\multicolumn{2}{l}{Within-seed gain}     & $+$38 & $+$13.2\,pts \\
\bottomrule
\end{tabular}
\end{table}

The full v12b config achieves \textbf{76.0\%}---a new best-ever at
$n=2.0$\,\AA\ by a large margin ($+$9.1\,pts over v12's 66.9\%,
$+$12.2\,pts over v11's 63.8\%).  While cross-seed comparison adds
$\pm$3.5\,pt uncertainty, the within-seed gain of $+$13.2\,pts over
baseline makes the improvement unambiguous.

At the sample level, the full v12b config rescues 56 samples vs.\
baseline while losing only 18 (3.1:1 ratio).  Of the 56 rescued,
32 were oscillating (57\%), 13 almost\_converged (23\%), 7 ghost\_modes
(13\%), and 4 slow\_convergence (7\%).  This demonstrates that the
probe+late mechanism addresses \emph{all} major failure classes, not
just oscillation.

\paragraph{Cross-config convergence landscape.}
Across all 5 v12b configs: 139/287 (48\%) converge in all 5, 36/287
(13\%) fail in all 5, and 251/287 (87\%) converge in at least one config.
The 36 permanently failing samples represent a hard core that no kick
variant resolves at 20k steps---these likely require 50k steps or
structural interventions.

\subsubsection{Updated Best-Ever Table}
\label{sec:v12b-best-ever}

\begin{table}[h]
\centering
\caption{Best strict convergence rate at each noise level across all runs
  and seeds ($N=287$).  Updated through v12b.}
\label{tab:best-ever-v12b}
\begin{tabular}{ccll}
\toprule
Noise & Best rate & Configuration & Run \\
\midrule
\multicolumn{4}{l}{\emph{10k steps}} \\
0.5\,\AA & 256/287 (89.2\%) & RFO+Schl+PLS       & v11 \\
1.0\,\AA & 230/287 (80.1\%) & RFO+PLS             & v11 \\
1.5\,\AA & 182/287 (63.4\%) & RFO+Schl+PLS        & v10c \\
2.0\,\AA & 128/287 (44.6\%) & GQT+PLS             & v10c \\
\midrule
\multicolumn{4}{l}{\emph{20k steps}} \\
0.5\,\AA & 272/287 (94.8\%) & RFO+Schl+PLS (20k)  & v11 \\
1.0\,\AA & 256/287 (89.2\%) & RFO+PLS (20k)        & v12 \\
1.5\,\AA & 230/287 (80.1\%) & RFO+PLS (20k)        & v12 \\
2.0\,\AA & \textbf{218/287 (76.0\%)} & RFO+PLS+adapt+probe+late (20k) & \textbf{v12b} \\
\midrule
\multicolumn{4}{l}{\emph{50k steps}} \\
0.5\,\AA & 273/287 (95.1\%) & RFO+PLS (50k)  & v12 \\
1.0\,\AA & 275/287 (95.8\%) & RFO+PLS (50k)  & v12 \\
1.5\,\AA & 266/287 (92.7\%) & RFO+PLS (50k)  & v12 \\
\bottomrule
\end{tabular}
\end{table}

\subsubsection{v12b Summary and Next Steps}
\label{sec:v12b-summary}

\textbf{Main conclusions:}
\begin{enumerate}[nosep]
  \item \textbf{Probe kicks are the key innovation}: $+$26 net converged
    over adaptive-only, with a 2.4:1 rescue-to-loss ratio.  Probing a
    distance ladder in both $\pm$ directions and accepting the first
    $n_\text{neg}$-reducing displacement prevents spectrally harmful
    kicks that destabilize borderline trajectories.
  \item \textbf{Late escape is clean and effective}: $+$11 net with a
    3.2:1 rescue-to-loss ratio.  Activating only after 15k steps avoids
    disrupting early-stage convergence.
  \item \textbf{Adaptive scaling alone is insufficient}: $+$1 net over
    fixed kicks.  Scaling magnitude to eigenvalue changes \emph{which}
    samples are rescued/lost but not the overall ratio.
  \item \textbf{Fixed kicks remain a wash}: both\_kicks without
    probe or adaptive shows 0 net gain in this seed, confirming v12's
    finding of high churn with fixed-scale perturbations.
  \item \textbf{76.0\% at $n=2.0$\,\AA\ is new best-ever} by
    $+$9.1\,pts over v12 and $+$12.2\,pts over v11 (at 20k steps).
  \item \textbf{69 failures remain}: dominated by oscillating (54/69,
    78\%).  Almost\_converged and ghost\_modes are largely resolved
    (6 and 4 remaining vs.\ 17 and 13 at baseline).
\end{enumerate}

\textbf{Next steps:}
\begin{itemize}[nosep]
  \item \textbf{50k + v12b kicks}: The two breakthroughs (50k steps from
    v12, probe+late kicks from v12b) have not yet been combined.  At
    $n=2.0$\,\AA, 50k steps should rescue additional oscillating failures
    beyond 20k, and probe+late kicks should accelerate convergence within
    the 50k budget.  This is the highest-priority experiment.
  \item \textbf{v12b at $n \leq 1.5$\,\AA}: The probe and late-escape
    mechanisms were only tested at $n=2.0$\,\AA.  Running at lower noise
    levels would test whether the improvement transfers or is
    $n=2.0$-specific.
  \item \textbf{Probe+late without adaptive}: Since adaptive scaling adds
    negligible value, a simpler config (both\_kicks+probe+late, without
    adaptive) might perform identically.  This simplification should be
    verified.
  \item \textbf{Late-escape timing}: The 15k activation threshold was
    chosen ad~hoc.  Testing 10k or 12k might capture additional samples
    that converge between 10k and 15k without the late mechanism.
  \item \textbf{Hard-core analysis}: The 36 samples that fail in all 5
    v12b configs need trajectory inspection.  Are they structurally
    trapped (persistent negative eigenvalues) or merely slow (would
    converge at 50k+)?
\end{itemize}

%=============================================================================
\section{What We Incorporated in v10/v10b}
\label{sec:incorporated}
%=============================================================================

\subsection{ARC Full-Spectrum Step Builder (v10) --- FAILED}

Replaces both the shifted Newton shift and trust region with cubic
regularization (Section~\ref{sec:arc}).  Implementation:
\texttt{\_nr\_step\_arc()} solves the secular
equation~\eqref{eq:arc-secular}; the ARC step-control loop
(lines~2350--2450 of \texttt{minimization.py}) manages $\sigma$-adaptation.

\textbf{Rationale}: The trust radius collapse mechanism
(Section~\ref{sec:collapse}) is the dominant failure mode.  ARC avoids
collapse by design---$\sigma$ can grow arbitrarily large (cautious) or
shrink arbitrarily small (aggressive), without a hard floor.

\textbf{Grid parameters} (v10 SLURM):
$\sigma_0 \in \{0.1, 1.0, 10.0\}$,
$\gamma_1 \in \{2.0, 3.0\}$.

\textbf{Result}: \textcolor{red}{\textbf{Catastrophic failure.}}  Best ARC
config: 23\% at $n=0.5$, 2\% at $n=1.0$, 0\% at $n\geq 1.5$.
80\% of all ARC trajectories classified as \texttt{energy\_plateau}:
$\sigma$ ratchets upward on noisy PES (even good steps can have poor
$\rho$), steps shrink to near-zero, energy stalls.
No $\sigma_0$/$\gamma_1$ combination recovers.  GDIIS does not rescue.
\textbf{ARC should not be used for noisy PES minimization.}

\subsection{GDIIS Oscillation Breaking (v10, corrected v10b) --- FAILED}

\texttt{\_GDIISExtrapolator} stores geometry--gradient--step history and
attempts interpolation every $N$ steps when oscillation is detected.

\textbf{v10b correction}: Changed error vectors from raw gradients to Newton
step vectors (Section~\ref{sec:gdiis}), per Schlegel's
formulation~\eqref{eq:gdiis-B}.

\textbf{Grid parameters}: buffer sizes $\{8, 12\}$.

\textbf{Result}: \textcolor{red}{\textbf{Zero effect.}}  Across all
configurations (buffer 8/12, oscillation-triggered, late-stage threshold
0.01), convergence counts and failure mode distributions are identical
with and without GDIIS.  Hypothesis: GDIIS extrapolation requires
error vectors to span a low-dimensional subspace; in $\sim$60-DOF
molecules, oscillation involves too many coupled modes for 8--12
history vectors to capture.

\subsection{RFO Step Builder (v10b)}

Augmented Hessian method (Section~\ref{sec:rfo}).
\texttt{\_solve\_rfo\_secular()} + \texttt{\_nr\_step\_rfo()}.
Uses the same trust-region step-control as v9 baseline.

\textbf{Rationale}: RFO is the standard in production codes (Gaussian, ORCA,
Sella).  It provides adaptive shift determination without the hyperparameter
$\varepsilon$.  Never tested in our codebase before.

\textbf{Result}: Without PLS, RFO underperforms GQT (67/48/32/12\% vs.\
75/55/44/27\%).  \textbf{With PLS, RFO becomes champion}
(89/78/63/42\%), a +22/+34/+32/+28\,pt improvement over bare RFO.
Interpretation: RFO provides better step \emph{directions} than GQT
but worse step \emph{lengths}; PLS fixes the length.

\subsection{Schlegel Trust Radius Fixes (v10b)}

Boundary-check growth~\eqref{eq:schlegel-grow} and step-anchored
shrink~\eqref{eq:schlegel-shrink}.  Gated by
\texttt{schlegel\_trust\_update=True}.

\textbf{Rationale}: The standard rules can shrink the radius below the
actual step size, wasting recovery.  The boundary check prevents
unnecessary growth.  Both are well-established in the
literature~\citep{schlegel2011}.

\textbf{Result}: Mixed.  Alone: $\pm$2--5\,pts depending on noise level
(slightly helps GQT at $n=1.0$, slightly hurts at $n=0.5$).
\textbf{Combines well with PLS}: RFO+Schlegel+PLS is champion.

\subsection{Late-Stage GDIIS Acceleration (v10b) --- FAILED}

Parameter \texttt{gdiis\_late\_force\_threshold}: triggers GDIIS
when \texttt{force\_norm} drops below threshold, regardless of oscillation
detection.

\textbf{Rationale}: Schlegel~\citep{schlegel2011} describes GDIIS as
accelerating final convergence routinely.  Population~A failures (blind
modes) have moderate forces but no oscillation---late-stage GDIIS could
help where oscillation-triggered GDIIS does not.

\textbf{Result}: Zero effect, same as oscillation-triggered GDIIS
(Section~\ref{sec:incorporated}, GDIIS subsection).

\subsection{Polynomial Line Search (v10c) --- BREAKTHROUGH}
\label{sec:pls}

Cubic interpolation using energy and gradient at current and previous
geometry.  Zero extra evaluations.  Fits $p(t) = at^3 + bt^2 + ct + d$
satisfying $p(0) = E_k$, $p'(0) = \nabla E_k \cdot d$,
$p(1) = E_{k+1}$, $p'(1) = \nabla E_{k+1} \cdot d$.  Steps to the
cubic minimizer if it lies in $[0, 1]$ and reduces energy.

\textbf{Rationale}: Schlegel~\citep{schlegel2011} describes cubic/quartic
energy fitting near inflection points.  We expected modest gains for
the eigenvalue-crossing regime.

\textbf{Result}: \textcolor{blue}{\textbf{The single most impactful
intervention across all versions.}}  +7 to +34 points depending on
configuration and noise level.  Halves the oscillating failure population.
Transforms RFO from underperformer (67/48/32/12\%) to champion
(89/78/63/42\%).

%=============================================================================
\section{What We Did NOT Incorporate and Why}
\label{sec:not-incorporated}
%=============================================================================

\subsection{Internal Coordinates}
\label{sec:internal-coords}

The literature unanimously agrees that internal coordinates (redundant
internals~\citep{pulay-redundant}, TRIC~\citep{wang-tric},
FCW~\citep{lindh-fcw}) improve Hessian conditioning for floppy modes.

\textbf{Why not}: Our molecules are small ($\leq 7$ heavy atoms, $\leq 23$
total).  With full analytical Hessians at every step, the conditioning
advantage is marginal.  The performance cliff at 2.0\,\AA\ is driven by
being far from quadratic, not by coordinate choice.  Converting to internal
coordinates would require a significant refactor of the Cartesian-only
codebase for likely marginal benefit.

\subsection{Sella's RS-RFO with Geodesic Back-Transformation}
\label{sec:sella}

Sella~\citep{sella} implements Restricted-Step RFO (RS-RFO) with iterative
diagonalization and geodesic back-transformation from internal to Cartesian
coordinates.

\textbf{Why not}: Requires an internal-coordinate framework.  Our codebase is
Cartesian-only.  The plain RFO we implemented captures the essential
augmented-Hessian mechanism without the coordinate transformation machinery.

\subsection{Polynomial Line Search Near Inflection Points}

\textbf{Status}: \textcolor{blue}{\textbf{Implemented in v10c}} and
proved to be the most impactful intervention.  See
Section~\ref{sec:pls}.  Originally deferred as future work in v10b;
implemented by another agent based on the v10b plan.

\subsection{Model-Hessian Preconditioning}

Packwood et al.~\citep{packwood2018} use a Lindh-type model Hessian as
preconditioner.  The key trick: drop gradient-dependent $H^{(2)}$ terms and
take $|d^2V/d\zeta^2|$ to guarantee positive definiteness.

\textbf{Why not}: We already have the full analytical Hessian.
Preconditioning is designed for approximate or absent Hessians.  Our problem
is direction quality at large displacement, not Hessian approximation error.

\subsection{Coordinate Freezing}

Schlegel (1989)~\citep{schlegel1989}: Freeze modes with negative eigenvalues,
converge remaining DOFs, then release.

\textbf{Why not}: Our negative eigenvalue directions rotate between steps
(overlap $< 0.5$ for ghost band, measured in v6 analysis).  Freezing a
specific direction when it changes every step is ineffective.

\subsection{SCINE SCF Convergence Investigation}

Tighter SCF convergence thresholds could potentially reduce the Hessian noise
floor (currently $\sim8\times10^{-3}$).

\textbf{Status}: Investigation script written (v10c) but crashed with
\texttt{NameError: name `defaultdict' is not defined}.  The bug is a
missing \texttt{from collections import defaultdict} import.
\emph{Still open as future work.}

\subsection{Transition1x Frequency Verification}

\textbf{Status}: \emph{Completed in v10c.}  DFTB0 frequency analysis on
all 287 DFT-labeled endpoints:
\begin{itemize}[nosep]
  \item Reactants: 29/287 (10.1\%) are genuine minima ($n_\text{neg}=0$)
    under DFTB0.
  \item Products: 30/287 (10.5\%) are genuine minima.
  \item Median force at ``minima'': 1.39\,eV/\AA\ (reactant), 1.16\,eV/\AA\
    (product)---$10{,}000\times$ above our convergence threshold.
  \item Correlation with optimizer failure: Pearson $r=0.22$, Spearman
    $r=0.17$.  Point-biserial (is\_minimum vs.\ failure): $r=-0.26$.
\end{itemize}
The weak correlations confirm that optimizer failure is \textbf{not primarily
driven by the starting-point PES character}.  The 90\% ``saddle rate'' at
DFT-labeled minima reflects genuine PES disagreement between DFT and DFTB0,
not a numerical artifact.

%=============================================================================
\section{Explicit Non-Recommendations}
\label{sec:non-recommendations}
%=============================================================================

The following have been \emph{tested and found harmful}.  Literature confirms
they should not be revisited.

\begin{enumerate}[nosep]
  \item \textbf{iHiSD crossover} (v8): Tested $\mu_\text{max} \in \{0, 0.1,
    0.5, 1.0, 2.0\}$.  Pure SN ($\mu=0$) best.  Theory: gradient alignment
    assumption violated at our displacements~\citep{ihisd}.

  \item \textbf{Line search} (v7): TR wins at all noise levels.  Theory:
    commits to poor direction~\citep{schlegel2011}.
    \emph{Note}: The polynomial line search (v10c) \emph{is} a line search,
    but within the trust region step---it refines step \emph{length}, not
    direction.  This distinction is critical.

  \item \textbf{LM damping} (v2): SN better.  Theory: over-regularizes
    positive eigenvalues.

  \item \textbf{Weight capping / SPDN} (v5): 0\% convergence.  Discards
    essential curvature information.

  \item \textbf{Stagnation escape}: Not helpful across versions.

  \item \textbf{v4 features}: Neg-trust-floor, blind-mode correction,
    aggressive trust recovery, bidirectional escape, mode-following.
    All tested, all hurt.

  \item \textbf{PIC-ARC} (first-order + cubic reg): 0\% convergence.
    First-order methods cannot navigate this PES~\citep{packwood2018}.

  \item \textbf{ARC} (v10): Catastrophic.  $\sigma$-adaptation fails on
    noisy PES---even good steps produce poor $\rho$, causing $\sigma$ to
    ratchet upward.  80\% energy plateau.  All $\sigma_0$/$\gamma_1$
    combinations tested.

  \item \textbf{GDIIS} (v10/v10b): Zero effect.  Tested oscillation-triggered
    (buffer 8, 12) and late-stage (force threshold 0.01).  Convergence and
    failure distributions identical with/without.
\end{enumerate}

%=============================================================================
\section{Grid Design and Results}
\label{sec:grid}
%=============================================================================

\subsection{v10/v10b Grid (48 configs)}

The v10 grid tested 12 configurations $\times$ 4 noise levels = 48
combinations, each with 287 samples.  See Table~\ref{tab:v10c-full}
for results.

\textbf{Key findings}:
\begin{itemize}[nosep]
  \item ARC: catastrophic across all $\sigma_0$, $\gamma_1$ (best 23\%).
  \item GDIIS: zero effect on all configurations.
  \item RFO: underperforms GQT without PLS.
  \item Schlegel TR: mixed $\pm$5\,pts.
\end{itemize}

\subsection{v10c Grid (+PLS, 56 configs)}

Added polynomial line search (PLS) to GQT and RFO+Schlegel configs:
14 configs $\times$ 4 noise = 56 total.

\textbf{Key findings}:
\begin{itemize}[nosep]
  \item PLS: +7 to +34\,pts across all configs.
  \item RFO+Schlegel+PLS: new champion (89/78/63/42\%).
  \item GQT+PLS: competitive, wins at 2.0\,\AA\ (45\% vs.\ 42\%).
\end{itemize}

\subsection{v11 Grid (All-PLS Factor Analysis, 32 configs)}

All 8 algorithms use PLS.  8 configs $\times$ 4 noise = 32 total.
Tests: (1)~GQT+Schlegel+PLS (never tested), (2)~RFO+PLS without Schlegel,
(3)~20k steps (RFO and GQT), (4)~trust floor 0.005\,\AA\ (RFO and GQT).

\textbf{Key findings} (full analysis in Section~\ref{sec:v11}):
\begin{itemize}[nosep]
  \item 20k steps: $+$6 to $+$21\,pts, primarily resolving oscillation.
  \item Trust floor 0.005: catastrophic ($-$10 to $-$30\,pts).
  \item Schlegel: neutral when PLS is active ($\pm$1.7\,pts).
  \item RFO+PLS: wins at all noise levels over GQT+PLS.
  \item \textbf{New best-ever}: 272/287 (94.8\%) at 0.5\,\AA, 183/287
    (63.8\%) at 2.0\,\AA\ (both RFO+Schl+PLS, 20k).
\end{itemize}

\subsection{v12 Grid (Kicks + 50k Steps, 17 configs)}

All configs use RFO+PLS.  Tests: (1)~oscillation kick, (2)~blind-mode
kick, (3)~both kicks, (4)~both kicks with eigvec1, (5)~50k step budget.
14 configs at $n=0.5$--1.5\,\AA\ (v12 main run) $+$ 3 configs at
$n=2.0$\,\AA\ (v12\_quick run).

\textbf{Key findings} (full analysis in Section~\ref{sec:v12}):
\begin{itemize}[nosep]
  \item 50k steps: $+$2 to $+$13\,pts over 20k, lossless (zero samples lost).
  \item Kicks at $n \leq 1.5$\,\AA: never fired, zero effect.
  \item Kicks at $n=2.0$\,\AA: $+$5.9\,pts net, high churn (32 rescued, 15 lost).
  \item Kick direction (eigvec0 vs.\ eigvec1) irrelevant.
  \item \textbf{New best-ever}: 275/287 (95.8\%) at 1.0\,\AA, 266/287
    (92.7\%) at 1.5\,\AA\ (both RFO+PLS, 50k).
\end{itemize}

\subsection{v12b Grid (Adaptive Probe Kicks + Late Escape, 5 configs)}

All configs at $n=2.0$\,\AA, 20k steps, RFO+PLS.  Tests cumulative
features: (1)~both\_kicks, (2)~+adaptive scaling, (3)~+probe kicks,
(4)~+late escape.  Single run (job~1146854), $N=287$.

\textbf{Key findings} (full analysis in Section~\ref{sec:v12b}):
\begin{itemize}[nosep]
  \item Probe kicks: $+$26 net converged (2.4:1 rescue-to-loss ratio).
  \item Late escape: $+$11 net converged (3.2:1 ratio).
  \item Adaptive scaling alone: $+$1 net (marginal).
  \item Fixed kicks: net zero (21 rescued, 21 lost).
  \item \textbf{New best-ever at $n=2.0$\,\AA}: 218/287 (76.0\%),
    $+$13.2\,pts over baseline within-seed.
\end{itemize}

\subsection{Recommended Configuration for Future Work}

Based on all evidence through v12b, the recommended baseline is:
\begin{center}
\fbox{\textbf{RFO + PLS + probe kicks + late escape, 50k steps}}
\end{center}
Schlegel trust radius updates are neutral when PLS is active
(Section~\ref{sec:v11-factors}) and can be omitted for simplicity.
The 50k step budget is justified at $n \geq 1.0$\,\AA\ ($+$6.6 to
$+$12.5\,pts over 20k, Table~\ref{tab:v12-50k}), with diminishing
returns at $n=0.5$ ($+$2.1\,pts).  The trust radius floor should remain
at 0.01\,\AA\ (lowering to 0.005\,\AA\ is catastrophic,
Table~\ref{tab:trf-effect}).

For throughput-constrained applications, RFO+PLS+probe+late at 20k steps
is the recommended budget-limited configuration.  At $n=2.0$\,\AA, this
achieves 76.0\% (Table~\ref{tab:v12b}), up from 62.7\% without kicks.
Adaptive kick scaling can be omitted without performance loss
(Section~\ref{sec:v12b-incremental}).

%=============================================================================
\section{Open Questions}
\label{sec:open}
%=============================================================================

\begin{enumerate}
  \item \textbf{Persistently hard samples}: In v11, samples 179, 199, 129
    are hardest (5--7/32 configs converge).  These share a common spectral
    signature: multiple near-zero negative eigenvalues with very low
    gradient overlap ($<0.1$).  \emph{Partially resolved:}
    Table~\ref{tab:hardest-10} reveals only three molecular formulae
    among the 10 hardest samples: C$_3$H$_5$NO$_2$, C$_5$H$_5$NO, and
    C$_3$H$_8$N$_2$O---all nitrogen+oxygen-containing organics.
    Trajectory analysis (Section~\ref{sec:trajectory-analysis}) shows
    permanently hard samples like~179 have strongly negative eigenvalues
    ($\lambda_\text{min} = -0.09$ at step 20k) that decrease monotonically
    but never reach zero, unlike rescued samples which cross zero at
    $\sim$13{,}500--17{,}500 steps.
    Structural analysis of the specific N--O bond rearrangement coordinates
    is still needed.

  \item \textbf{Why does PLS help RFO so much more than GQT?}
    RFO+Schlegel+PLS gains +22--34\,pts over bare RFO, while GQT+PLS gains
    only +7--18\,pts.  Hypothesis: RFO's augmented Hessian gives better
    \emph{directions} but worse \emph{step lengths}; PLS fixes the length.
    v11 confirms: RFO+PLS (no Schlegel) outperforms GQT+PLS by 3--14\,pts
    at all noise levels, showing the RFO step direction is genuinely
    superior.

  \item \textbf{Population~A (blind modes) largely resolved by v12b}:
    Almost\_converged and ghost\_mode failures were previously unaffected
    by any intervention including 20k steps.  v12b's probe kicks and
    late escape reduce these from 17+13=30 to 6+4=10 at $n=2.0$\,\AA.
    \emph{Partially resolved:} the remaining 10 blind-mode failures
    (and the 54 oscillating failures) now constitute the residual
    failure population.  The 36 samples that fail across all 5 v12b
    configs likely require 50k steps or structural interventions.

  \item \textbf{GQT vs.\ RFO crossover at 2.0\,\AA}: \emph{Partially
    resolved.}  v11 shows RFO+PLS wins at all noise levels in this seed
    ($+$3.1\,pts at 2.0\,\AA).  The v10c crossover (GQT winning at
    2.0\,\AA) was seed-dependent and within the $\pm$3.5\,pt
    uncertainty.  No noise-adaptive selection is needed; RFO+PLS
    is weakly but consistently preferred.

  \item \textbf{SCINE SCF convergence}: Could tighter SCF thresholds reduce
    the Hessian noise floor ($\sim8\times10^{-3}$)?  Investigation script
    crashed (\texttt{NameError}: \texttt{defaultdict} not imported).
    \emph{Still open.}

  \item \textbf{Transition1x frequency verification}: Completed in v10c.
    Only 10\% of DFT-labeled minima are true minima under DFTB0.
    Weak correlation ($r=0.22$) with optimizer failure.
    \emph{Resolved}: not a primary driver of failure.

  \item \textbf{Optimal step budget}: v12 confirmed 50k steps gives
    $+$2 to $+$12.5\,pts over 20k (Section~\ref{sec:v12-50k}).
    \emph{Partially resolved.}  The remaining question is whether
    50k steps \emph{combined with} v12b's probe+late kicks (not yet
    tested) would push $n=2.0$\,\AA\ convergence above 80\%.
    The 36 samples that fail across all 5 v12b configs at 20k are
    the primary targets for this combined experiment.

  \item \textbf{Why does lower trust floor hurt?}  v11 shows floor
    0.005\,\AA\ is catastrophic ($-$10 to $-$30\,pts).  The
    $\varepsilon^{1/2}$ theorem predicts floor $\sim$0.01; the
    experimental confirmation strengthens this connection.  However, the
    mechanism deserves further theoretical analysis: is the floor effect
    purely about step size, or does it interact with the trust-radius
    recovery dynamics?
\end{enumerate}

%=============================================================================
\section{Codebase Architecture}
\label{sec:architecture}
%=============================================================================

\subsection{Pipeline}

\begin{verbatim}
run_minimization_parallel.py   (CLI + SLURM interface)
  |-> ParallelSCINEProcessor  (spawns N workers)
        |-> scine_worker_sample()  (creates DFTB0 predict_fn)
              |-> run_newton_raphson()  (core optimizer, ~2700 lines)
                    |-> predict_fn(coords, atomic_nums, do_hessian=True)
                    |   -> (energy, forces, hessian)
                    |-> get_vib_evals_evecs()  (Eckart decomposition)
                    |-> Step builder:
                    |     _nr_step_arc()       (v10: ARC secular eq)
                    |     _nr_step_rfo()       (v10b: RFO augmented Hessian)
                    |     _nr_step_shifted_newton()  (v3-v9: GQT)
                    |     _nr_step_lm_damping()      (v2, inferior)
                    |     _nr_step_hard_filter()      (v1, inferior)
                    |-> Step control:
                    |     ARC sigma-adaptation  (v10, FAILED)
                    |     Trust region           (v1-v9, RFO)
                    |     Polynomial line search (v10c, BEST)
                    |     Line search            (v7, inferior)
                    |-> GDIIS:
                    |     _GDIISExtrapolator    (v10/v10b, ZERO EFFECT)
                    |     Oscillation detection  (v10, no benefit)
                    |     Late-stage trigger      (v10b, no benefit)
                    |-> Convergence:
                          Cascade evaluation at 8 thresholds
                          Strict / Relaxed / Stalled classification
\end{verbatim}

\subsection{Key Files}
\begin{itemize}[nosep]
  \item \texttt{baselines/minimization.py} --- Core optimizer ($\sim$2700 lines)
  \item \texttt{runners/run\_minimization\_parallel.py} --- CLI and parallelization
  \item \texttt{scripts/analyze\_minimization\_nr\_grid.py} --- Grid analysis, tag parsing
  \item \texttt{scripts/analyze\_nr\_failure\_autopsy.py} --- Failure classification
  \item \texttt{scripts/diagnostic\_from\_minima.py} --- DFTB0 forensics at DFT minima
  \item \texttt{scripts/slurm\_templates/} --- SLURM job scripts (v3--v10)
\end{itemize}

%=============================================================================
\section*{References}
%=============================================================================

\begin{thebibliography}{99}

\bibitem[Cartis et al.(2011a)]{arc-part1}
  C.~Cartis, N.~I.~M.~Gould, Ph.~L.~Toint,
  ``Adaptive cubic regularisation methods for unconstrained optimization.
  Part~I: Motivation, convergence and numerical results,''
  \emph{Math.\ Program.}, 127(2):245--295, 2011.

\bibitem[Cartis et al.(2011b)]{cgt-optimal}
  C.~Cartis, N.~I.~M.~Gould, Ph.~L.~Toint,
  ``Optimal Newton-type methods for nonconvex smooth optimization problems,''
  Technical Report, 2011.

\bibitem[Cartis et al.(2012)]{cgt-complexity}
  C.~Cartis, N.~I.~M.~Gould, Ph.~L.~Toint,
  ``Complexity bounds for second-order optimality in unconstrained optimization,''
  \emph{J.~Complexity}, 28(1):93--108, 2012.

\bibitem[Boumal et al.(2018)]{boumal-arc-manifolds}
  N.~Boumal, P.-A.~Absil, C.~Cartis,
  ``Global rates of convergence for nonconvex optimization on manifolds,''
  \emph{IMA J.\ Numer.\ Anal.}, 39(1):1--33, 2019.

\bibitem[Benson and Shanno(2014)]{benson-shanno}
  H.~Y.~Benson, D.~F.~Shanno,
  ``Interior-point methods for nonconvex nonlinear programming:
  cubic regularization,''
  \emph{Comput.\ Optim.\ Appl.}, 58(2):323--346, 2014.

\bibitem[Schlegel(2011)]{schlegel2011}
  H.~B.~Schlegel,
  ``Geometry optimization,''
  \emph{WIREs Comput.\ Mol.\ Sci.}, 1(5):790--809, 2011.

\bibitem[Schlegel(1989)]{schlegel1989}
  H.~B.~Schlegel,
  ``A comparison of geometry optimization with internal, Cartesian,
  and mixed coordinates,''
  \emph{Int.\ J.\ Quantum Chem.}, 44(S23):243--252, 1993.

\bibitem[Pulay(1980)]{pulay-diis}
  P.~Pulay,
  ``Convergence acceleration of iterative sequences.\ The case of SCF iteration,''
  \emph{Chem.\ Phys.\ Lett.}, 73(2):393--398, 1980.

\bibitem[Cs\'asz\'ar and Pulay(1984)]{gdiis-csaszar-pulay}
  P.~Cs\'asz\'ar, P.~Pulay,
  ``Geometry optimization by direct inversion in the iterative subspace,''
  \emph{J.\ Mol.\ Struct.}, 114:31--34, 1984.

\bibitem[Li and Frisch(2006)]{li-frisch-gediis}
  X.~Li, M.~J.~Frisch,
  ``Energy-represented direct inversion in the iterative subspace within a
  hybrid geometry optimization method,''
  \emph{J.\ Chem.\ Theory Comput.}, 2(3):835--839, 2006.

\bibitem[Pulay and Fogarasi(1992)]{pulay-redundant}
  P.~Pulay, G.~Fogarasi,
  ``Geometry optimization in redundant internal coordinates,''
  \emph{J.\ Chem.\ Phys.}, 96(4):2856--2860, 1992.

\bibitem[Wang and Song(2016)]{wang-tric}
  L.-P.~Wang, C.~Song,
  ``Geometry optimization made simple with translation and rotation
  internal coordinates,''
  \emph{J.\ Chem.\ Phys.}, 144(21):214108, 2016.

\bibitem[Lindh et al.(1999)]{lindh-fcw}
  R.~Lindh, A.~Bernhardsson, M.~Sch\"utz,
  ``Force-constant weighted internal coordinates in vibrational frequency
  calculations,''
  \emph{Chem.\ Phys.\ Lett.}, 303(5--6):567--575, 1999.

\bibitem[Packwood et al.(2018)]{packwood2018}
  D.~M.~Packwood, J.~Kermode, L.~Mones, N.~Bernstein, J.~R.~Kermode,
  G.~Cs\'anyi, C.~Ortner,
  ``A universal preconditioner for simulating condensed phase materials,''
  \emph{Sci.\ Rep.}, 8:5374, 2018.

\bibitem[Su et al.(2022)]{ihisd}
  Y.~Su, L.~Zhang, W.~E,
  ``The iterative high-index saddle dynamics method for the index-$k$
  saddle point on the energy landscape,''
  arXiv:2602.01883, 2022.

\bibitem[Luo et al.(2023)]{convergence-degenerate}
  L.~Luo, Y.~Nesterov, H.~He, E.~A.~Gryazina,
  ``Convergence of regularized Newton methods on degenerate critical points,''
  2023.

\bibitem[Hermes et al.(2022)]{sella}
  E.~D.~Hermes, K.~Sargsyan, H.~N.~Najm, B.~C.~Garrett,
  ``Sella: a Python package for finding transition states on
  potential energy surfaces,''
  \emph{J.\ Chem.\ Theory Comput.}, 18(11):6974--6988, 2022.

\bibitem[Bhowmik et al.(2023)]{transition1x}
  A.~Bhowmik, I.~Batatia, S.~Batzner, et al.,
  ``Transition1x: A dataset for building generalizable reactive machine
  learning potentials,''
  \emph{Sci.\ Data}, 10:232, 2023.

\bibitem[ORCA manual]{orca-manual}
  F.~Neese et al., ``ORCA Manual,'' version~5.0, Max-Planck-Institut
  f\"ur Kohlenforschung, 2023.

\end{thebibliography}

\end{document}
